% !TeX program = pdflatex
% ─────────────────────────────────────────────────────────────────────────────
%  psi/multifocal.tex — A INTELIGÊNCIA MULTIFOCAL, FUNDADA.
%
%  O QUE ESTE FICHEIRO É: a teoria da inteligência multifocal de Augusto Jorge
%  Cury («Inteligência Multifocal — Análise da Construção dos Pensamentos e da
%  Formação de Pensadores», Cultrix, 8.ª ed.) posta em forma: postulados
%  declarados, definições, teoremas e demonstrações, sobre o alicerce do
%  `fisica.tex`.
%
%  DUAS PARTES, e a ordem importa.
%   I — A CONSTRUÇÃO. Seis postulados de tradução, e tudo o resto sai deles com
%       demonstração: a reconstrução, o eu como acumulação, a multifocalidade,
%       a âncora, a psicoadaptação, os dois pensamentos conscientes, a
%       interpretação do outro, o custo do caos e a medida do «antiespaço».
%   II — AS RACHADURAS. Catorze sítios onde o livro afirma duas coisas que não
%       podem valer juntas, e o teorema da Parte I que fecha cada uma.
%
%  A REGRA: a rachadura é sempre um PAR de asserções DO PRÓPRIO LIVRO. Não se
%  discute com o autor a partir de fora. Onde ele está certo, diz-se que está —
%  e a §pe mostra que é na maior parte.
%
%  E NÃO SE TOCA NO fisica.tex. Este documento CITA-O pelo rótulo, e repete o
%  enunciado citado em uma linha, para que isto se leia sem abrir outro
%  ficheiro. A matéria formal do alicerce é de lá e fica lá; o que se demonstra
%  aqui é a ESPECIALIZAÇÃO dela aos postulados desta tradução.
%
%  Compila SOZINHO (`pdflatex psi/multifocal.tex`) e dentro do livro.
% ─────────────────────────────────────────────────────────────────────────────
\documentclass[../livro.tex]{subfiles}

\begin{document}

% ─── Licença ──────────────────────────────────────────────────────────────
% Proprietário / comercial. Copyright (c) 2026 Aarão Melo Lopes.
% Uso de código ou textos mediante contrato e pagamento — ver LICENSE na raiz.
% Contacto: aarao.melo.lopes@gmail.com

% ── o que o estilo.tex não dá, e este documento usa ──
% O teste é `\lema` e NÃO `\c@lema`: o lema partilha o contador `teorema`, pelo
% que `\c@lema` nunca chega a existir e o guarda seria sempre falso.
\ifcsname lema\endcsname\else
  \theoremstyle{plain}\newtheorem{lema}[teorema]{Lema}
\fi
\providecolor{cinza}{gray}{0.35}
\providecommand{\abs}[1]{\lvert #1\rvert}

%% A CITAÇÃO DO LIVRO. É sempre literal, e vai entre aspas angulares.
\providecommand{\cury}[1]{\emph{«#1»}}

%% A CITAÇÃO DO fisica.tex: o nome do resultado, e o rótulo por baixo, para que
%% a dívida se possa cobrar sem sair desta página.
\providecommand{\fis}[2]{\textbf{#2}~\textnormal{(\code{#1})}}

%% O POSTULADO DE TRADUÇÃO, citado pelo número: são seis, e são todos.
\providecommand{\pt}[1]{\textbf{P#1}}

%% A LINHA DE FECHO de cada rachadura.
\providecommand{\sepc}{\quad\penalty0$\big|$\quad\penalty0}
\providecommand{\fecha}[3]{%
\par\smallskip\noindent{\footnotesize\raggedright
\textbf{no livro}: #1\sepc\textbf{a rachadura}: #2\sepc\textbf{o fecho}: #3\par}\smallskip}

\ifSubfilesClassLoaded{%
  \gkcapa{A Inteligência Multifocal, Fundada}
         {seis postulados, doze teoremas, catorze rachaduras}
         {a teoria de Augusto Cury posta em forma sobre a dobra --- o que se
          demonstra, o que não fechava, e o que fica por fechar}
  \maketitle
  \tableofcontents
}{%
  \part*{A Inteligência Multifocal}\addcontentsline{toc}{part}{A Inteligência Multifocal}
}

\begin{abstract}\noindent
\textbf{O que se leu.} \emph{Inteligência Multifocal --- Análise da Construção
dos Pensamentos e da Formação de Pensadores}, de Augusto Jorge Cury (Editora
Cultrix, 8.\textsuperscript{a} edição revista e ampliada, 2006; primeira edição
1998). O livro inteiro, do prefácio ao glossário.

\medskip\noindent\textbf{O que o livro afirma.} Que a memória não guarda para
devolver: \cury{não existe lembrança}, só reconstrução. Que a leitura da memória
é feita por \textbf{quatro} fenómenos e não por um --- a autochecagem, o
autofluxo, a âncora da memória e o \emph{eu} ---, e que três deles são
inconscientes, pelo que a maior parte do que se pensa não foi decidido por
ninguém. Que a interpretação tem \textbf{cinco etapas}, das quais as três
primeiras são inconscientes, e que o \emph{eu} \textbf{só se forma na quinta}.
Que a emoção se adapta ao que se repete --- a \textbf{psicoadaptação} --- e que
por isso nem a dor nem o prazer duram. E que entre a consciência e a coisa há um
\cury{antiespaço}, uma distância que nunca se atravessa.

\medskip\noindent\textbf{O que este documento faz --- e são duas coisas, nesta
ordem.} \emph{Primeiro constrói.} Declara \textbf{seis postulados de tradução} e
mais nenhum, e deduz deles, com demonstração, as afirmações centrais da teoria:
que a leitura é reconstrução e não resgate e \emph{quanto} se perde; que o
\emph{eu} é uma \textbf{acumulação} e os outros três fenómenos são o
\textbf{incremento}; que nenhuma leitura carrega o seu autor --- e é isso, e não
a contagem quatro, que faz a inteligência ser multifocal; que a âncora é uma
\textbf{bola} e que os territórios particionam; que a psicoadaptação é
$1/G$; que a dor se perpetua por repartição e se alivia por concentração; que os
dois pensamentos conscientes são as duas metades de um operador, e a cisão é
única; que a traição da interpretação é a \textbf{soma}; e que o antiespaço tem
medida e ela halva. \emph{Depois diagnostica.} Encontra no livro
\textbf{catorze} pares de afirmações do próprio autor que não podem valer
juntas, e fecha cada um com um teorema da primeira parte.

\medskip\noindent\textbf{E o resultado, numa linha.} Quase tudo o que o livro
descreve \emph{como fenómeno} é, nesta construção, \textbf{uma leitura de um
campo}; e quase todas as contradições do livro nascem de confundir duas leituras
diferentes \textbf{do mesmo campo} com dois fenómenos diferentes:
\[
\boxed{\;\text{a disponibilidade é }\zeta;\qquad
\text{a resposta é }\mu;\qquad
\text{a adaptação é }1/G\;}
\]
\emph{Três nomes, um campo.} O livro tinha razão em dizer que são inevitáveis, e
não tinha como dizer porquê --- porque lhe faltava a conta.

\medskip\noindent\textbf{E o que fica de pé é mais do que o que cai.} As teses
centrais de Cury --- que a inteligência é multifocal, que o \emph{eu} não é o
autor da maior parte do que pensa, que a leitura da memória é reconstrução, e que
nenhuma classificação de pessoas se lê no funcionamento da mente --- saem daqui
\textbf{demonstradas} em vez de observadas. O que cai são as junções, e são as
junções que este documento refaz.
\end{abstract}

% ═══════════════════════════════════════════════════════════════════════════
\section{A regra deste documento, e o que ele não afirma}\label{mf:regra}

\noindent\textbf{Quatro regras, e são de leitura obrigatória} --- sem elas este
documento lê-se como uma coisa que não é.

\medskip\noindent\textbf{(1) A construção parte de seis postulados, e de mais
nenhum.} Os postulados da §\ref{mf:postulados} são a \emph{totalidade} do
compromisso com a psicologia. Tudo o que a Parte~I afirma depois deles é
demonstrado --- ou de um resultado do \code{fisica.tex}, citado pelo rótulo, ou
dos postulados, ou dos dois. \emph{Quem negar um postulado sabe exactamente que
teoremas caem com ele}, porque cada demonstração diz quais usou.

\medskip\noindent\textbf{(2) A rachadura é interna, e é sempre um par.} Não se
chama rachadura a uma tese com que se discorde, nem a uma afirmação sem prova ---
um livro de teoria tem direito a postulados. Chama-se rachadura ao sítio onde
\textbf{o livro afirma $A$ e afirma $\neg A$}, ou onde afirma $A$ e afirma $B$ e
$B$ torna $A$ impossível. As duas asserções vão citadas \emph{à letra}, e o
capítulo delas vai ao lado. \emph{Quem discordar da rachadura vai ao livro e
confere as duas frases.}

\medskip\noindent\textbf{(3) O fecho é um teorema desta obra, e não uma
sugestão.} Cada rachadura da Parte~II fecha num enunciado da Parte~I, com o
número. \emph{Nenhum fecho é inventado no sítio}: se o resultado não existisse, a
rachadura ficaria na lista dos que não fecham (§\ref{mf:aberto}), e há-os.

\medskip\noindent\textbf{(4) E o fecho custa um compromisso, que se declara.}
Fechar uma rachadura é escolher uma leitura, e a leitura escolhida exclui outras.
Onde isso acontece --- e acontece em todas --- diz-se \textbf{qual é o
compromisso}. \emph{Um fecho que não custe nada não fechou nada.}

\medskip
\gkcaixa{\textbf{O que este documento NÃO afirma.} \emph{Nada aqui é uma
afirmação sobre cérebros.} O \code{fisica.tex} constrói um operador, uma
realização e um campo; não constrói neurónios, não mede pessoas, e não tem uma
única observação clínica. Quando o Teor.~\ref{mf:thm:adapta} diz que \emph{a
psicoadaptação é $1/G$}, afirma \textbf{uma coisa e só uma}: que dos postulados
\pt{1}--\pt{3} isso se segue. \textbf{Não se afirma} que a mente humana
\emph{seja} uma realização $\pi\colon I\to X$; não se afirma que os números
medidos nos programas da bateria digam alguma coisa sobre um paciente; e não se
afirma que Cury estivesse a descrever isto. \emph{Os medidos citados são medidas
de programas, e de mais nada.} Quem quiser passar daqui para a clínica tem de
medir na clínica --- e este documento não o fez. A §\ref{mf:aberto} diz o que
isso deixa em aberto, e diz-o pelo nome.}

% ═══════════════════════════════════════════════════════════════════════════
% ═══════════════════════════════════════════════════════════════════════════
\part{A construção}

\section{O alicerce, e o que dele se importa}\label{mf:alicerce}

\noindent\textbf{Este documento não reconstrói a álgebra}, e diz de onde a
tira. Tudo o que segue assenta em cinco objectos do \code{fisica.tex}, e são
estes --- \emph{cada um com o seu rótulo, para que a dívida se possa cobrar}:

\medskip
{\small
\begin{tabular}{@{}lll@{}}
\toprule
o objecto & o que é & onde \\
\midrule
$\partial$ & a \textbf{dobra}: um mapa com $\partial\circ\partial=\mathrm{id}$ & \code{fis:def:op} \\
$\pi\colon I\to X$ & a \textbf{realização}: $X$ só precisa de igualdade decidível & \code{fis:def:objeto} \\
$G(x)=\abs{\pi^{-1}(x)}$ & a \textbf{multiplicidade}: quantos índices na célula $x$ & \code{fis:def:objeto} \\
$\Phi=\sum_x(G(x)-1)$ & a \textbf{folga}: quantos vieram parar a sítio ocupado & \code{fis:def:folga} \\
$\zeta,\ \mu=\zeta^{-1}$ & \textbf{acumular} e \textbf{desacumular}: $\mu$ é a diferença finita & \code{fis:thm:mu} \\
\bottomrule
\end{tabular}}

\medskip\noindent\textbf{E uma frase do alicerce, que é a que este documento
mais usa.} A multiplicidade obedece a uma recorrência local em que
\emph{nenhum passo consulta $\pi(0),\dots,\pi(t)$}: cada escrita altera $G$ num
único ponto, e esse ponto é a célula corrente
(\fis{fis:thm:mult}{Teor. «multiplicidade geométrica»}, cláusula~(3)). É a isso
que a obra chama \emph{a memória não pertence ao agente} --- e é dela que sai
quase tudo o que segue.

% ───────────────────────────────────────────────────────────────────────────
\section{Os seis postulados de tradução}\label{mf:postulados}

\noindent\textbf{São seis, e são todos.} Cada um traduz um objecto do livro num
objecto do alicerce. \emph{Não são teoremas e não se demonstram} --- são o preço
de entrada, e é por isso que estão juntos, numerados, e numa página só.

\begin{quote}\small
\textbf{\pt{1} (o registo).} A história intrapsíquica é uma realização
$\pi\colon I\to X$, onde $I$ é a lista \emph{ordenada} das experiências
vividas e $X$ o conjunto das representações psicossemânticas, com igualdade
decidível e mais nada. \emph{O fenómeno RAM \textbf{é} $\pi$}: escreve, e não
pergunta.

\smallskip\textbf{\pt{2} (a disponibilidade).} A disponibilidade histórica da
representação $x$ no instante $t$ é $G_t(x)$.

\smallskip\textbf{\pt{3} (a leitura).} O que um passo \emph{lê} não é o campo: é
a \textbf{diferença relativa} do campo, $G(y)/G(x)-1$, sobre a vizinhança
corrente.

\smallskip\textbf{\pt{4} (o território).} $X$ é hierárquico: as representações
endereçam-se por palavras de $p$ bits, com a distância
$d(a,b)=2^{-\operatorname{prof}(a,b)}$ da árvore do encaixe
(\code{fis:def:arvore}). \emph{A âncora da memória é a bola $\mathcal{B}(a,r)$
dessa árvore.}

\smallskip\textbf{\pt{5} (o consciente).} O \emph{eu} no instante $t$ é a
\textbf{acumulação} do campo até $t$; os fenómenos da autochecagem, do autofluxo
e da âncora são os \textbf{incrementos} que a compõem.

\smallskip\textbf{\pt{6} (o outro).} A leitura de outra pessoa é a
\textbf{soma}, numa matriz única, do padrão diretivo e dos padrões associativos
que o estímulo convoca.
\end{quote}

\medskip\noindent\textbf{E diga-se logo qual é o mais caro.} É \pt{4}: pedir que
a memória seja hierárquica é uma hipótese forte, e é a mais forte deste
documento --- as três propriedades da âncora saem dela e caem com ela. \pt{3} é
o segundo: identificar «o que se sente» com o gradiente, e não com o campo, é
uma escolha, e é a que faz a psicoadaptação fechar. \pt{1}, \pt{2} e \pt{5} são
baratos: o livro já os afirma quase nas mesmas palavras. \pt{6} é a regra de
Hebb, e usa-se \emph{só} na §\ref{mf:outro}.

\medskip\noindent\textbf{E o que \emph{não} está aqui.} Não há postulado para
\emph{energia psíquica}. O livro põe-na como substância; este documento não a
traduz e não a usa, e a §\ref{mf:aberto} diz porquê.

% ───────────────────────────────────────────────────────────────────────────
\section{A memória: o que se conserva, o que se perde, e quanto}\label{mf:memoria}

\begin{teorema}[a memória conserva a contagem e perde a fibra]\label{mf:thm:reconstroi}
Sob \pt{1}--\pt{2}:
\begin{enumerate}\itemsep3pt
\item[(1)] \textbf{Conserva:} $\displaystyle\sum_{x}G(x)=\abs{I}$, seja qual for
a história.
\item[(2)] \textbf{E não distingue:} existem duas histórias $\pi\neq\pi'$ sobre
o mesmo $I$ com a mesma medida e dobras diferentes, isto é, com
$\max G_\pi\neq\max G_{\pi'}$.
\item[(3)] \textbf{Logo a leitura não é resgate.} A aplicação
$\pi\mapsto\bigl(\sum_xG(x)\bigr)$ não é injetiva; o estado da memória não
determina a história que o produziu.
\end{enumerate}
\end{teorema}

\begin{proof}
(1) é o \fis{fis:lem:conserva}{Lema «conservação»}: as fibras $\pi^{-1}(x)$
particionam $I$, e a cardinalidade é aditiva sobre partições finitas.

(2) é o \fis{fis:thm:medida}{Teor. «a medida conserva-se, a fibra perde-se»}, e a
testemunha está dentro da própria definição: sobre o mesmo $I$ com
$\abs{I}\ge2$ e um $X$ com pelo menos dois objectos, tome-se $\pi$
\textbf{constante} --- donde $\max G=\abs{I}$ --- e $\pi'$ \textbf{não
constante} --- donde $\max G'<\abs{I}$. As duas cumprem (1).

(3) é (2) lido como enunciado sobre a aplicação: duas pré-imagens distintas com a
mesma imagem.
\end{proof}

\medskip\noindent\textbf{É isto, e só isto, que \cury{não existe lembrança}
quer dizer que se possa demonstrar.} \emph{Não} que nada se conserve --- a
cláusula~(1) diz que a contagem se conserva sempre. É que \textbf{o que se
conserva não chega para reconstituir}: a marca não diz quantas patas a
produziram.

\begin{corolario}[o débito tem número]\label{mf:cor:debito}
Sob \pt{1}--\pt{2}, com $\pi$ sobrejetiva sobre $X$, o que se perdeu é
\[
\boxed{\;\Phi\;=\;\sum_{x}\bigl(G(x)-1\bigr)\;=\;\abs{I}-\abs{X}\;}
\]
e $\Phi=0$ \emph{exactamente} quando nenhuma experiência caiu sobre outra.
\end{corolario}

\begin{proof}
É a \fis{fis:def:folga}{Def. «a folga, local e total»} com o
\fis{fis:cor:folga}{Cor. «a folga é uma só»}: a soma da esquerda é
$\sum_xG(x)-\abs{\operatorname{supp}G}$, que por (1) do
Teor.~\ref{mf:thm:reconstroi} vale $\abs{I}-\abs{X}$; e ela anula-se exactamente
quando $G\le1$, que é a injetividade
(\fis{fis:thm:mult}{Teor. «multiplicidade geométrica»}, cláusula~(2)).
\end{proof}

\medskip\noindent\textbf{E daqui sai a leitura fina da recordação.} O que fica de
uma experiência não é a experiência: é um par. Na secção
\code{fis:rastro} a obra chama-lhe o \textbf{rastro}, e separa três coisas com
três nomes ---
$(m,\delta)$ é o par, $m$ é o que se guarda, e $\delta$ é a perda \emph{medida}.
A frase da casa é a que aqui interessa: \emph{$\delta$ não é um erro a lamentar;
é o que a dobra devolve, e é por existir que se sabe quanto ficou de fora}.

\begin{corolario}[duas recordações da mesma coisa são duas folhas]\label{mf:cor:folhas}
Sob \pt{1}, para cada representação $x$ os números de visita percorrem
\emph{exactamente} $\{1,2,\dots,G(x)\}$, e o par $\bigl(\pi(i),k(i)\bigr)$ é
injetivo. Logo duas recordações da mesma experiência \textbf{não são a mesma
leitura repetida}: são dois pontos distintos da fibra vertical de $x$.
\end{corolario}

\begin{proof}
É o \fis{fis:thm:levant}{Teor. «levantamento em folhas»}, cláusulas~(3) e~(1):
fixado $x$ e sendo $i_1<\dots<i_g$ a sua fibra, $k(i_r)=r$, que é uma bijeção
sobre $\{1,\dots,g\}$; e daí a injetividade de $\tilde\pi=(\pi,k)$.
\end{proof}

\medido{\code{tests/aranha\_n.c} §AN1--§AN2: campo incremental igual à recontagem
e $\sum G=\abs{I}$. §AN16: \emph{a mesma medida com dobras diferentes} --- as
duas curvas, com a reta como controlo. \code{tests/meio.c} §M8: o centro contra
as suas fronteiras em três grades, $\mathbf{6}$ a $\mathbf{13}$ vezes melhor.
\emph{São medidas de programas.}}

% ───────────────────────────────────────────────────────────────────────────
\section{O \emph{eu} é uma acumulação, e gerir é desacumular}\label{mf:eu}

\noindent A peça que o livro não tem é a \textbf{diferença de nível} entre o
\emph{eu} e os outros três fenómenos. Ela existe, é uma identidade, e traz com
ela o atraso do \emph{eu} --- com o seu tamanho exacto.

\begin{teorema}[o campo escreve com $\zeta$ e lê-se com $\mu$]\label{mf:thm:zetamu}
Sob \pt{1}, \pt{2} e \pt{5}, pondo
$a_x(t)=\mathbf{1}_{\{\pi(t)=x\}}$:
\begin{enumerate}\itemsep3pt
\item[(1)] \textbf{Os três mordomos são o incremento:} cada passo altera $G$
numa célula só, e $a_x$ é esse incremento.
\item[(2)] \textbf{O \emph{eu} é a acumulação:} $G_t(x)=(a_x*\zeta)(t)$.
\item[(3)] \textbf{Gerir é desacumular:} $a_x(t)=G_t(x)-G_{t-1}(x)$; conhecida a
família $\{G_t\}$, a trajetória inteira volta --- \emph{sem que nada dela tenha
sido guardado}.
\item[(4)] \textbf{E o atraso é exactamente um passo.} $G_t$ depende de $a_x(u)$
apenas para $u\le t$; logo o \emph{eu} em $t-1$ não contém $a_x(t)$, e o
\emph{eu} em $t$ é o \emph{eu} em $t-1$ mais $a_x(t)$.
\end{enumerate}
\end{teorema}

\begin{proof}
(1) é a cláusula~(3) do \fis{fis:thm:mult}{Teor. «multiplicidade geométrica»}:
$G_{t+1}(x)=G_t(x)+\mathbf{1}_{\{x=\pi(t+1)\}}$, e para $x\neq\pi(t+1)$ a
indicatriz anula-se.

(2) e (3) são as cláusulas~(1) e~(2) do
\fis{fis:thm:zetamu}{Teor. «o ciclo é $\zeta$; a recuperação é $\mu$»}, com
$\mu=\zeta^{-1}$ a ser a diferença finita
(\fis{fis:thm:mu}{Teor. «$\mu=\zeta^{-1}$ na ordem total»}).

(4) De (2), $G_t(x)=\sum_{u\le t}a_x(u)$, soma em que não entra nenhum $u>t$; e a
diferença $G_t-G_{t-1}$ é o único termo que os separa, que é $a_x(t)$ por (3).
\end{proof}

\begin{corolario}[porque é mais fácil ser vítima do que agente]\label{mf:cor:vitima}
Sob \pt{5}, \textbf{uma acumulação não pode ser um incremento}: o \emph{eu} em
$t$ é função de $a$ até $t$, e o passo em $t$ não é função do \emph{eu} em $t$
--- é uma das suas parcelas. Logo o \emph{eu} não determina o passo que o
constitui, e determina-o com um passo de atraso e não mais.
\end{corolario}

\begin{proof}
A primeira metade é a cláusula~(4) do Teor.~\ref{mf:thm:zetamu}. Para a segunda:
por (3) do mesmo teorema, $a_x(t)$ recupera-se de $\{G_t,G_{t-1}\}$, que estão
ambos disponíveis em $t$; logo o atraso é de um passo e não é maior.
\end{proof}

\medskip\noindent\textbf{E é isto que dá sentido ao \cury{retroativamente} do
livro.} Não é agir no passado --- o que não quer dizer nada. É \textbf{ler o
passado do campo presente}: $\mu$ opera sobre $\{G_t\}$, que está todo escrito
agora, e devolve o passo que foi dado. \emph{A operação clínica que o livro
descreve e não sabe nomear é uma deconvolução.}

\medido{\code{lib/incidencia.h} com \code{tests/zetamu.c}:
$\zeta\circ\mu=\mu\circ\zeta=\mathrm{id}$ com resíduo $\mathbf{0}$ em toda a
largura, e a trajetória $a_x$ recuperada de $\{G_t(x)\}$ exactamente nas visitas.
\textbf{8 unidades, zero falhas.} \code{tests/aranha\_n.c} §AN15:
$\mu=\zeta^{-1}$ em $\mathbf{4096}$ pares.}

% ───────────────────────────────────────────────────────────────────────────
\section{A multifocalidade, demonstrada --- e não é a contagem quatro}\label{mf:multi}

\noindent A tese que dá nome ao livro pode enunciar-se e demonstrar-se, e o
enunciado \textbf{não menciona quantos são os leitores}. É melhor assim: a
contagem quatro é da observação do autor, e o que faz o trabalho é outra coisa.

\begin{teorema}[nenhuma leitura carrega o seu autor]\label{mf:thm:multifocal}
Sob \pt{1}--\pt{2}, atribua-se a cada escrita um agente. Então:
\begin{enumerate}\itemsep3pt
\item[(1)] \textbf{Invariância:} duas execuções que produzam a mesma ocupação
por célula dão campos iguais célula a célula, e $\abs{I}$ é o mesmo --- $G$ é
invariante sob qualquer reatribuição de agentes que preserve a ocupação.
\item[(2)] \textbf{Logo nenhuma função do campo distingue autores.} Se $F$ é
qualquer aplicação definida em $G$, então $F$ toma o mesmo valor nas duas
execuções.
\item[(3)] \textbf{E a indistinção é a dos dois extremos:} um agente a visitar
uma célula $K$ vezes e $K$ agentes a visitá-la uma vez cada produzem o mesmo
campo.
\end{enumerate}
\end{teorema}

\begin{proof}
(1) e (3) são o \fis{fis:thm:agentes}{Teor. «$G$ não vê o construtor»}: pela
cláusula~(3) do \fis{fis:thm:mult}{Teor. «multiplicidade geométrica»}, cada
escrita incrementa $G$ na célula corrente e \emph{nada mais regista} --- em
particular nenhuma escrita carrega a identidade de quem a fez; como as duas
execuções produzem a mesma multiplicidade por célula, os campos coincidem.

(2) é imediato de (1): uma função de um invariante é um invariante. Se
$G=G'$ então $F(G)=F(G')$.
\end{proof}

\medskip\noindent\textbf{E aqui está a tese do livro, dita como ela é.} A
inteligência não é multifocal \emph{porque há quatro fenómenos}; é multifocal
porque \textbf{o campo não tem coluna para o autor}. Acrescentar ou tirar
leitores não muda nada no ciclo, e é isso que faz o mecanismo escalar:
\[
\boxed{\;\text{o campo é memória da \textbf{construção}, e não do
\textbf{construtor}}\;}
\]

\begin{teorema}[a introspecção repõe a ordem, e não a autoria]\label{mf:thm:folha}
Sob \pt{1}, seja $k(i)$ o número de visita e $\tilde\pi(i)=(\pi(i),k(i))$. Então:
\begin{enumerate}\itemsep3pt
\item[(1)] $\tilde\pi$ é \textbf{injetiva}, $\tilde G=1$ na imagem, e $G$
recupera-se de $\tilde\pi$ sem qualquer outro dado. \emph{Uma coordenada, e a
dobra desfaz-se} --- e ela sobe \emph{um} andar, seja qual for a dimensão.
\item[(2)] \textbf{Mas o agente não é função de $\tilde\pi$}: existem duas
atribuições de agentes distintas com o mesmo $\tilde\pi$. Separar agentes pede
outra coordenada --- o identificador --- e essa \textbf{não sai do campo}.
\end{enumerate}
\end{teorema}

\begin{proof}
(1) é o \fis{fis:thm:levant}{Teor. «levantamento em folhas»}, cláusulas~(1),
(3) e~(5).

(2) A atribuição de agentes é um dado \emph{a mais} sobre a mesma trajetória:
fixada a sucessão $\pi$ --- e portanto fixados $k$ e $\tilde\pi$ --- pode
repartir-se as escritas por um agente ou por vários sem alterar $\pi(t)$ para
nenhum $t$. Logo $\tilde\pi$ é o mesmo e a atribuição não; donde ela não é
função de $\tilde\pi$. Que a distinção também não sai de $G$ é a cláusula~(2) do
Teor.~\ref{mf:thm:multifocal}.
\end{proof}

\medskip\noindent\textbf{Donde a leitura clínica, e ela é mais modesta do que a
do livro e mais útil.} O \cury{stop introspectivo} não é --- e não pode ser ---
uma auditoria de autoria: ninguém recupera \emph{quem} construiu um pensamento.
O que a introspecção repõe é a \textbf{ordem}: o que veio primeiro, o que veio a
seguir, e quantas vezes se passou por ali. \emph{É o que os clínicos de facto
fazem}, e agora sabe-se porquê.

\medido{\code{tests/aranha\_n.c} §AN31: um agente com $4$ voltas e $4$ agentes
com uma volta --- campos iguais nas $\mathbf{64}$ células, $G=4$ nas $16$, com o
controlo negativo a ver as $16$ a diferir contra $3$ voltas. E a separação medida
dos dois lados: a ordem separa \textbf{todos} os pares que caem na mesma célula
($\mathbf{96}$ de $96$), o identificador separa só os que atravessam a fronteira
entre agentes ($\mathbf{64}$).}

% ───────────────────────────────────────────────────────────────────────────
\section{A âncora é a bola, e os territórios particionam}\label{mf:ancora}

\noindent O livro declara a âncora \cury{difícil de ser definida} e faz dela a
peça que suporta o edifício clínico. \pt{4} dá-lhe a definição, e as três
propriedades que o livro precisava de postular passam a ser cláusulas.

\begin{teorema}[o território]\label{mf:thm:ancora}
Sob \pt{4}, com $\mathcal{B}(a,r)=\{x:d(a,x)\le r\}$ e $r\ge0$:
\begin{enumerate}\itemsep3pt
\item[(1)] \textbf{Delimita, e é uma célula:} $\pi_r(i)=$ a bola de raio $r$ que
contém $i$ é uma realização no sentido de \code{fis:def:objeto}; a sua célula
\emph{é} a bola, e $G$ conta os índices dentro dela.
\item[(2)] \textbf{É uma:} duas bolas de raio $r$ ou coincidem ou são disjuntas;
para cada $r$ as bolas \textbf{particionam} $X$. Há \emph{exactamente} uma que
contém o ponto corrente.
\item[(3)] \textbf{Alarga-se de dentro:} todo ponto é centro, e para
$r'\ge r$ tem-se $\mathcal{B}(a,r)\subseteq\mathcal{B}(a,r')$ com o \emph{mesmo}
ponto por centro. \emph{Subir o raio não pede um ponto de fora.}
\item[(4)] \textbf{E há dois deslocamentos, e não um:} mudar $r$ (o território
fica mais largo ou mais estreito) e mudar de prefixo (o território muda de sítio
à mesma profundidade). \emph{São operações diferentes.}
\end{enumerate}
\end{teorema}

\begin{proof}
(1) e (2) são as cláusulas~(3) e~(2) do
\fis{fis:thm:bolas}{Teor. «as bolas particionam»}, que valem porque $d$ é
ultramétrica e \textbf{simétrica}
(\fis{fis:lem:ultra}{Lema «$d$ é ultramétrica, e é simétrica»}).

(3) A primeira metade é a cláusula~(1) do mesmo teorema: se
$b\in\mathcal{B}(a,r)$ então $\mathcal{B}(b,r)=\mathcal{B}(a,r)$. A inclusão é
imediata: $d(a,x)\le r\le r'$. E o centro conserva-se porque
$a\in\mathcal{B}(a,r')$ e vale outra vez a cláusula~(1).

(4) As bolas identificam-se pelo prefixo (cláusula~(3) do mesmo teorema, na
árvore finita de $p$ bits): um par $(\text{prefixo},r)$, e cada coordenada varia
sem a outra.
\end{proof}

\begin{corolario}[não há duas leituras ao mesmo tempo]\label{mf:cor:umsotem}
Sob \pt{4}, fixado $r$, o ponto corrente pertence a uma e uma só bola de raio
$r$. \emph{É a versão exacta de \cury{não há convivência simultânea de dois
tipos de emoções, motivações e pensamentos na mente}.}
\end{corolario}

\begin{proof}
Cláusula~(2) do Teor.~\ref{mf:thm:ancora}: as bolas de raio $r$ formam uma
partição, e um elemento pertence a exactamente uma parte de uma partição.
\end{proof}

\begin{corolario}[a proximidade entre territórios tem número]\label{mf:cor:prox}
Sob \pt{4}, para dois padrões de $N$ níveis que concordam nos primeiros $q$ e
discordam nos seguintes, a sobreposição é $(2q-N)/N$, \emph{estritamente
crescente} em $q$. Duas memórias próximas são duas memórias na mesma bola.
\end{corolario}

\begin{proof}
É o \fis{fis:thm:sobrepoe}{Teor. «a sobreposição é a contagem dos níveis que
concordam»}: há $q$ concordâncias e $N-q$ discordâncias, e a derivada em $q$ é
$2/N>0$. A ligação à bola é a distância da árvore, $d=2^{-\operatorname{prof}}$,
que é decrescente na mesma profundidade.
\end{proof}

\medido{\code{tests/aranha\_n.c} §AN6: as bolas particionam; todo ponto é
centro; \emph{e a euclidiana falha ambas com o mesmo raio}.
\code{tests/hopfield.c} §F3: a sobreposição é exactamente a conta dos níveis que
concordam, em $\mathbf{4096}$ pares sem um erro, e cresce com o prefixo.}

% ───────────────────────────────────────────────────────────────────────────
\section{A psicoadaptação é $1/G$, e o luto crónico é subaditividade}\label{mf:adapta}

\noindent Aqui a construção dá mais do que o livro pedia: dá a lei, e dá com ela
o mecanismo de duas coisas que o livro observa e não explica --- a perpetuação da
dor e o alívio pela exposição.

\begin{teorema}[a adaptação é a divisão pela contagem]\label{mf:thm:adapta}
Sob \pt{1}--\pt{3}, seja $x$ uma representação com disponibilidade $G\ge1$ e
considere-se mais uma visita a $x$. Então:
\begin{enumerate}\itemsep3pt
\item[(1)] \textbf{O que se escreve é sempre o mesmo:} o incremento vale $1$,
seja qual for $G$.
\item[(2)] \textbf{O que se lê decresce:} a diferença relativa entre o campo
depois e antes é
\[
\boxed{\;\frac{G+1}{G}-1=\frac{1}{G}\;}
\]
\item[(3)] \textbf{E é estritamente decrescente em $G$:} $1/(G+1)<1/G$ para todo
$G\ge1$. Logo a resposta a repetições sucessivas do mesmo estímulo decresce
\emph{enquanto o registo cresce}.
\item[(4)] \textbf{E não depende da escala:} trocando $G$ por $cG$ com $c>0$, a
diferença relativa não se move.
\end{enumerate}
\end{teorema}

\begin{proof}
(1) é a cláusula~(1) do Teor.~\ref{mf:thm:zetamu}: cada escrita incrementa $G$ na
célula corrente em uma unidade.

(2) é \pt{3} aplicado ao par (antes, depois): a diferença relativa é
$(G+1)/G-1$, e a conta é imediata. Que a grandeza lida seja esta, e não a
diferença absoluta, é a cláusula~(2) do
\fis{fis:thm:gradiente}{Teor. «o gradiente é a diferença relativa»} --- ao
levantar o índice com $g^{uv}=G^{-1}h^{uv}$ a escala cancela e o que resta é a
razão menos um.

(3) $1/(G+1)<1/G$ porque $G<G+1$ e $t\mapsto1/t$ é decrescente em $t>0$.

(4) é a mesma cláusula~(2): o diferencial multiplica por $c$ e o levantamento
divide por $c$.
\end{proof}

\medskip\noindent\textbf{E as duas asserções contrárias do livro passam a ser a
mesma frase, lida em dois sítios:}

\begin{center}\small
\begin{tabular}{@{}lll@{}}
\toprule
no livro & é & e por isso \\
\midrule
o RAM regista em zonas privilegiadas & $G$ sobe & fica \textbf{mais} disponível \\
a emoção adapta-se ao repetido & $1/G$ desce & sente-se \textbf{menos} \\
\bottomrule
\end{tabular}
\end{center}

\noindent\emph{Uma frase é sobre o campo, a outra é sobre o gradiente do campo, e
a segunda é o inverso da primeira.}

\begin{lema}[a soma harmónica é subaditiva]\label{mf:lem:sub}
Seja $H_g=\sum_{r=1}^{g}1/r$, com $H_0=0$. Então, para $a,b\ge0$ inteiros,
\[
H_{a+b}\;\le\;H_a+H_b ,
\]
com igualdade se e só se $a=0$ ou $b=0$. Além disso $H_g\le g$, com igualdade se
e só se $g\le1$.
\end{lema}

\begin{proof}
$H_{a+b}-H_a=\sum_{r=a+1}^{a+b}1/r$, e cada termo desta soma é
$1/(a+r')\le1/r'$ para $r'=1,\dots,b$, com desigualdade \emph{estrita} sempre que
$a\ge1$. Logo $H_{a+b}-H_a\le H_b$, estritamente se $a\ge1$ e $b\ge1$. A segunda
afirmação é a mesma conta com $1/r\le1$, estrita para $r\ge2$.
\end{proof}

\begin{corolario}[o luto crónico é repartição; a exposição é concentração]\label{mf:cor:luto}
Sob \pt{1}--\pt{3}, seja uma história de $N$ visitas repartida por $k$
representações distintas com fibras de tamanhos $g_1,\dots,g_k$,
$\sum_ig_i=N$. Chame-se \textbf{resposta total} a
$R=\sum_{i}H_{g_i}$, a soma do que se leu em cada visita. Então:
\[
\boxed{\;H_N\;\le\;R\;\le\;N\;}
\]
--- o \textbf{mínimo} $H_N$ atinge-se \emph{concentrando} tudo numa
representação ($k=1$), e o \textbf{máximo} $N$ atinge-se \emph{espalhando} por
$N$ representações distintas ($k=N$, todos $g_i=1$).
\end{corolario}

\begin{proof}
Que a resposta acumulada numa fibra de tamanho $g$ é $H_g$: pela cláusula~(2) do
Teor.~\ref{mf:thm:adapta}, a $r$-ésima visita lê $1/r$, e a fibra tem os números
de visita $1,\dots,g$ exactamente uma vez (Cor.~\ref{mf:cor:folhas}); somar dá
$H_g$.

O limite superior: pelo Lema~\ref{mf:lem:sub}, $H_{g_i}\le g_i$, e somar dá
$R\le\sum_ig_i=N$, com igualdade sse todos os $g_i\le1$.

O limite inferior: por subaditividade e indução em $k$,
$H_N=H_{g_1+\dots+g_k}\le\sum_iH_{g_i}=R$, com igualdade sse todos menos um dos
$g_i$ forem nulos.
\end{proof}

\medskip\noindent\textbf{E é este o mecanismo que faltava ao livro, em duas
frases.} Cury observa que \cury{a multiplicidade das experiências psíquicas\dots
dificulta a operacionalização espontânea psicodinâmica do fenômeno da
psicoadaptação, estabelecendo uma perpetuação da dor, um luto crônico}.
\emph{Aqui isso é aritmética}: com o mesmo número $N$ de visitas, repartir
maximiza a resposta e concentrar minimiza-a, e a diferença entre os dois extremos
é $N$ contra $H_N$ --- que cresce como o logaritmo. \textbf{Não se sofre mais:
sofre-se em mais lados}, e é isso que não adapta. E a recíproca é a exposição:
voltar deliberadamente à mesma representação faz $G$ subir e $1/G$ cair.
\emph{É a mesma conta, corrida ao contrário.}

% ───────────────────────────────────────────────────────────────────────────
\section{Os dois pensamentos conscientes: o que mede e o que ordena}\label{mf:conscientes}

\noindent O livro tem três tipos de pensamento e, para os ligar, dois fenómenos
postulados --- o «LVD» e um segundo que ele deixa em \cury{provavelmente}. A
cisão certa é outra, é \textbf{única}, e dispensa os dois.

\begin{teorema}[a partição é única, e cada metade faz uma coisa só]\label{mf:thm:duas}
\emph{Hipótese, e é a mais forte desta parte}: suponha-se que a leitura do campo
sobre si mesmo se representa por uma matriz $W$ de pesos, com transposta, e que
a actualização é assíncrona por sinal. Então $W$ decompõe-se de maneira
\textbf{única} em $W=W_s+W_a$ com $W_s$ simétrica e $W_a$ antissimétrica, e:
\begin{enumerate}\itemsep3pt
\item[(1)] \textbf{$W_s$ mede:} a energia $E=-\tfrac12\sum_{ij}W_{ij}s_is_j$
nunca sobe, e a trajetória atinge um \emph{ponto fixo} em número finito de
passos.
\item[(2)] \textbf{$W_a$ não mede:} $s^{\!\top}W_a\,s=0$ para todo $s$, e a
dinâmica deixa de ter a função de Lyapunov que a energia fornecia.
\item[(3)] \textbf{E é só isto que a antissimetria dá.} O que a órbita faz depois
--- ciclar, e com que período --- \textbf{não se segue daqui}.
\end{enumerate}
\end{teorema}

\begin{proof}
Existência e unicidade da cisão, e as cláusulas~(1) e~(2), são o
\fis{fis:thm:duastorres}{Teor. «a partição é única, e cada metade faz uma coisa
só»}: a existência é
$W=\tfrac12(W+W^{\!\top})+\tfrac12(W-W^{\!\top})$, a unicidade sai de transpor e
somar; (1) é o \fis{fis:thm:desce}{Teor. «a descida, e porque ela não sobe»},
onde $\Delta E\le0$ e o espaço de estados é finito; e (2) é o cancelamento aos
pares, $(W_a)_{ji}=-(W_a)_{ij}$ com diagonal nula. A cláusula~(3) é a
advertência do mesmo teorema.
\end{proof}

\medskip\noindent\textbf{E é aqui que os dois pensamentos conscientes do livro
caem no sítio:}
\[
\boxed{
\begin{array}{lll}
\text{o \textbf{essencial}} & \text{o campo }G & \text{real, e no \textbf{espaço} --- não no agente}\\[2pt]
\text{o \textbf{dialético}} & W_s & \text{\emph{mede}: desce, tem ponto fixo, verbaliza-se}\\[2pt]
\text{o \textbf{antidialético}} & W_a & s^{\!\top}W_as=0:\ \text{\emph{define o indefinível}}
\end{array}}
\]

\begin{corolario}[a ansiedade vital é a metade que a energia não vê]\label{mf:cor:ansiedade}
Sob a hipótese do Teor.~\ref{mf:thm:duas}, a metade antissimétrica é
\emph{exactamente} a parte do operador que não aparece na energia e move a
órbita. Logo ela \textbf{não tem para onde parar}, e não por escolha: não há
grandeza que ela faça descer.
\end{corolario}

\begin{proof}
É a \fis{fis:prop:cego}{Prop. «a antissimetria é exactamente o que a energia não
vê e a dinâmica vê»}, junta com a cláusula~(2) do Teor.~\ref{mf:thm:duas}: sem
forma quadrática não há função de Lyapunov, e o argumento de terminação da
cláusula~(1) não tem sobre o que operar.
\end{proof}

\begin{corolario}[a passagem entre o campo e a sua leitura não custa]\label{mf:cor:ponte}
A travessia entre a célula e o espectro conserva a norma e é involutiva; logo
\textbf{não há custo por explicar} na passagem do campo aos pensamentos
conscientes, e não é preciso postular um fenómeno para a fazer.
\end{corolario}

\begin{proof}
É o \fis{fis:thm:ponte}{Teor. «a travessia entre a célula e o espectro não paga»}
com o \fis{fis:thm:H}{Teor. «inversa, involutividade e Parseval»}.
\end{proof}

\medskip\noindent\textbf{E a assimetria que o livro observa --- que o dialético
traduz mal o antidialético e não o contrário --- lê-se na transformada:}
\emph{o módulo não carrega nada, e a fase carrega tudo}, e a fase, em largura um,
é \textbf{um bit}. Quem lê pela face que mede recebe magnitude e perde o lado;
quem lê pela outra recebe o lado e não tem magnitude. \emph{O «truncamento da
comunicação» não é uma deficiência da linguagem: é a face pela qual se leu.}

\medido{\code{tests/hopfield.c} §F1: a energia nunca sobe num passo assíncrono,
em milhares de passos, e $8$ padrões corrompidos a $10\%$ voltam todos ao
original. §F8: a partição $W=W_s+W_a$ é exacta, medida \emph{em inteiros e por
igualdade}, sem limiar. §F10: \emph{falha} --- e a falha é o resultado, porque a
cláusula~(3) não afirma o ciclo.}

% ───────────────────────────────────────────────────────────────────────────
\section{A interpretação do outro: a traição é a soma}\label{mf:outro}

\noindent É a única secção que usa \pt{6}, e dela sai o mecanismo que o livro
não tem para uma palavra que ele usa muito: \cury{inevitável}.

\begin{teorema}[somar identifica; separar custa espaço]\label{mf:thm:outro}
Sob \pt{6}, com $W_{ij}=\sum_{p}\xi^p_i\xi^p_j$ a matriz onde os padrões
diretivo e associativos são somados:
\begin{enumerate}\itemsep3pt
\item[(1)] \textbf{A leitura não é injetiva:} conjuntos distintos de padrões
podem dar a mesma $W$; a multiplicidade dessa realização é $G>1$, e a isso
chama-se \textbf{interferência}.
\item[(2)] \textbf{E a razão é o núcleo:} um morfismo que soma tem núcleo, e
dois conjuntos cuja diferença esteja no núcleo têm a mesma imagem.
\item[(3)] \textbf{Desfazer custa uma coordenada:} pondo cada padrão no seu
ramo, a realização é injetiva e $G\le1$ --- é o levantamento, com o ramo no
lugar da folha.
\item[(4)] \textbf{E a coordenada ocupa lugar:} a matriz somada é $N^2$
\textbf{fixo} e satura; a árvore \textbf{cresce} com o que guarda. \emph{A folga
é uma só: quem enche paga com dobra, quem conta paga com buraco.}
\end{enumerate}
\end{teorema}

\begin{proof}
(1) e (3) são o \fis{fis:thm:interfere}{Teor. «somar identifica; separar não»}.
(2) é o \fis{fis:thm:palavra}{Teor. «a palavra reconstrói, e a dobra é o
núcleo»}: dois prefixos com a mesma imagem dão um factor no núcleo, o que é
$G>1$ pela cláusula~(2) do \fis{fis:thm:mult}{Teor. «multiplicidade
geométrica»}; na árvore cada memória ocupa um prefixo distinto e prefixos
distintos são bolas disjuntas (Teor.~\ref{mf:thm:ancora}(2)). (4) é o
\fis{fis:cor:folga}{Cor. «a folga é uma só»}, lido nos dois mapas.
\end{proof}

\begin{corolario}[conhecemos o outro a partir de nós, e isso é uma divisão]\label{mf:cor:outroeu}
Sob \pt{1}--\pt{3} e \pt{6}, a imagem do outro entra como incremento sobre um
campo que já carrega a história do observador; pelo Teor.~\ref{mf:thm:adapta}, o
que dela se lê é $1/G$ com $G$ a disponibilidade prévia. \emph{Quanto mais
história, menos o outro aparece} --- e o \cury{quanto} é $1/G$.
\end{corolario}

\begin{proof}
Cláusula~(2) do Teor.~\ref{mf:thm:adapta}, aplicada à célula onde o padrão do
outro é somado.
\end{proof}

\medskip\noindent\textbf{E a prescrição do livro fica com o seu preço}, que é a
parte que ele não tinha:
\[
\boxed{\;\text{a árvore troca \textbf{interferência} por \textbf{espaço}}\;}
\]
\emph{Colocar-se no lugar do outro é dar-lhe um ramo} --- não o somar à própria
matriz. E isso \textbf{ocupa lugar}: não é capacidade infinita de graça, e cada
distinção preservada ocupa lugar na hierarquia. Eis porque é que quase ninguém o
faz, e a resposta não é falta de virtude: \emph{é que a alternativa barata ---
somar --- é a que o mecanismo dá de borla}.

\medido{\code{tests/hopfield.c} §F5: a árvore nunca fica abaixo da matriz somada
e fica acima em pelo menos um $P$, \emph{sem limiar}; e o preço em espaço está
medido, não escondido. §F2: a saturação da matriz somada é o resultado ---
\emph{e o valor da constante não}.}

% ───────────────────────────────────────────────────────────────────────────
\section{O custo do caos, e a medida do «antiespaço»}\label{mf:custo}

\noindent Faltam duas contas, e são as que o livro pede sem as abrir: quanto
custa apagar, e qual é a distância entre a consciência e a coisa.

\begin{teorema}[o que custa é apagar, e a conta é a folga]\label{mf:thm:custo}
\emph{Hipótese importada e declarada}: apagar um bit tem um custo mínimo
$\kappa>0$. Sob \pt{1} e com $\pi$ sobrejetiva:
\begin{enumerate}\itemsep3pt
\item[(1)] \textbf{Destrutivo paga:} uma operação que identifica estados paga
por cada bit que perde.
\item[(2)] \textbf{Injetivo não paga:} o piso de uma operação que nunca leva dois
estados ao mesmo é \textbf{zero exacto} --- e não um número pequeno.
\item[(3)] \textbf{E a conta total é a folga:}
$\displaystyle\kappa\sum_x\bigl(G(x)-1\bigr)=\kappa\bigl(\abs{I}-\abs{X}\bigr)
=\kappa\,\Phi$, que se anula exactamente quando $\pi$ é injetiva.
\item[(4)] \textbf{Donde o dual guardado é o piso não pago:} guardar a folha do
Teor.~\ref{mf:thm:folha} põe o piso a zero, e paga em \emph{espaço} o que
deixaria de pagar em \emph{calor}.
\end{enumerate}
\end{teorema}

\begin{proof}
(1), (2) e (4) são o \fis{fis:thm:landauer}{Teor. «o piso do invertível é zero
exacto»}: perder um bit é ter dois estados com a mesma saída, que é $G>1$; ser
injetiva é $G\le1$, e sem fibra não trivial a soma dos custos é uma soma vazia;
e a reversão obtém-se guardando a coordenada que a dobra ia deitar fora, que é
$\tilde\pi=(\pi,k)$.

(3) é o \fis{fis:cor:custofolga}{Cor. «a conta da dissipação é a folga, e a
folga já estava contada»}, com o Cor.~\ref{mf:cor:debito} a dar a segunda
igualdade.
\end{proof}

\medskip\noindent\textbf{E é isto que põe a prescrição central do livro em
contas.} Cury defende que se \cury{revise, recicle, reorganize} em vez de se
descartar, e apresenta-o como higiene intelectual. Aqui é uma factura: \emph{o
que se guarda reverte a custo zero; o que se deita fora é cobrado}. E o critério
não é uma propriedade a inspeccionar --- é uma conta a fazer:
\emph{resíduo zero $\Longleftrightarrow$ reverte $\Longleftrightarrow$ piso
zero}.

\begin{teorema}[o antiespaço tem medida, e ela halva]\label{mf:thm:rastro}
Batendo alternadamente as duas faces sobre um par $(L,U)$ que cerca um ponto:
\begin{enumerate}\itemsep3pt
\item[(1)] \textbf{O que fica é um par:} $(m,\delta)$ com $m=\tfrac{L+U}{2}$ e
$\delta=\tfrac{U-L}{2}$ --- \emph{$(L,U)$ escrito duas vezes, com resíduo zero}.
\item[(2)] \textbf{A compressão é ficar só com $m$}, e o que ela deita fora não é
desconhecido: é $\delta$, \emph{medido}.
\item[(3)] \textbf{Há dois desfechos, e os dois fecham:} ou o ponto está na
grade, e o processo acaba em $m=g$; ou não está, e o processo acaba com a largura
no grão --- \emph{as duas faces cercam-no}.
\item[(4)] \textbf{E a largura desce:} a largura seguinte não passa de
$\delta^2/m$, e a cota é logarítmica.
\end{enumerate}
\end{teorema}

\begin{proof}
(1) e (2) são a secção \code{fis:rastro}, onde os três nomes --- rastro,
compressão, perda --- são separados. (3) é a dicotomia da mesma secção, com o
caso na grade a ser a cláusula~(3.I) do \fis{fis:thm:corte}{Teor. «o corte»}.
(4) é a cláusula~(2) do mesmo teorema, $(m-g)(m+g)=\delta^2$ com $m+g\ge m$, e o
\fis{fis:cor:cota}{Cor. «a cota é logarítmica»}.
\end{proof}

\medskip\noindent\textbf{Donde o \cury{antiespaço} do livro deixa de ser um
absoluto e passa a ser uma grandeza.} A consciência \textbf{não incorpora} a
coisa --- verdade: fica com $m$, e $m$ não é o ponto. E \textbf{umas leituras
estão mais perto do que outras} --- verdade: é $\delta$, e $\delta$ é um número
que desce. \emph{O que é infinito é o número de batidas que seriam precisas para
lá chegar, e não a distância a que se está.}

\begin{corolario}[a tríade do autor é a alternância das faces]\label{mf:cor:triade}
O encaixe da cláusula~(4) do Teor.~\ref{mf:thm:rastro} exige que se bata
\textbf{alternadamente} nas duas faces: batendo sempre do mesmo lado o intervalo
não encaixa. \emph{A arte da pergunta, a arte da dúvida e a arte da crítica são a
alternância}, e é essa a sua razão estrutural.
\end{corolario}

\begin{proof}
A cláusula~(1) do \fis{fis:thm:corte}{Teor. «o corte»}: é a aplicação alternada
das duas médias que dá $g\le g'\le m'\le m$, e é dela que sai a contracção da
cláusula~(2). Uma face sozinha não produz o encaixe.
\end{proof}

\medido{\code{tests/dissipa.c} --- $4$ blocos, $\mathbf{0}$ falhas. Na mesma
tarefa e com a mesma entrada, a destrutiva paga
$\mathbf{376\,304\,828\,416}$ rJ e a involutiva paga $\mathbf{0}$; e a
involutiva, repetida, devolve o original em $\mathbf{4096}$ de $4096$ posições.
\code{tests/meio.c} §M8: os dois desfechos medidos lado a lado --- o processo a
acabar em $m=g$ com o ponto na grade, e a torre de duplicação a acabar com o
ponto \emph{cercado}.}

% ───────────────────────────────────────────────────────────────────────────
\section{O que a Parte I derivou, e de quê}\label{mf:derivou}

{\small
\begin{longtable}{@{}p{0.34\textwidth}p{0.28\textwidth}p{0.30\textwidth}@{}}
\toprule
o que ficou demonstrado & de que postulados & sobre que resultado \\
\midrule
\endfirsthead
\toprule
o que ficou demonstrado & de que postulados & sobre que resultado \\
\midrule
\endhead
\bottomrule
\multicolumn{3}{@{}r@{}}{\footnotesize\itshape continua na página seguinte}\\
\endfoot
\bottomrule
\endlastfoot
Teor.~\ref{mf:thm:reconstroi} --- conserva a contagem, perde a fibra &
\pt{1}, \pt{2} & \code{fis:lem:conserva}, \code{fis:thm:medida} \\
Cor.~\ref{mf:cor:debito} --- o débito é $\Phi=\abs{I}-\abs{X}$ &
\pt{1}, \pt{2} & \code{fis:def:folga}, \code{fis:cor:folga} \\
Cor.~\ref{mf:cor:folhas} --- duas recordações são duas folhas &
\pt{1} & \code{fis:thm:levant} \\
Teor.~\ref{mf:thm:zetamu} --- o \emph{eu} é $\zeta$, gerir é $\mu$ &
\pt{1}, \pt{2}, \pt{5} & \code{fis:thm:zetamu}, \code{fis:thm:mu} \\
Cor.~\ref{mf:cor:vitima} --- o atraso é exactamente um passo &
\pt{5} & Teor.~\ref{mf:thm:zetamu}(4) \\
Teor.~\ref{mf:thm:multifocal} --- nenhuma leitura carrega o autor &
\pt{1}, \pt{2} & \code{fis:thm:agentes} \\
Teor.~\ref{mf:thm:folha} --- repõe-se a ordem, não a autoria &
\pt{1} & \code{fis:thm:levant} \\
Teor.~\ref{mf:thm:ancora} --- a âncora é a bola, e particiona &
\pt{4} & \code{fis:thm:bolas}, \code{fis:lem:ultra} \\
Cor.~\ref{mf:cor:umsotem} --- um só território de cada vez &
\pt{4} & Teor.~\ref{mf:thm:ancora}(2) \\
Cor.~\ref{mf:cor:prox} --- a proximidade tem número &
\pt{4} & \code{fis:thm:sobrepoe} \\
Teor.~\ref{mf:thm:adapta} --- a psicoadaptação é $1/G$ &
\pt{1}--\pt{3} & \code{fis:thm:gradiente}, \code{fis:thm:zetamu} \\
Cor.~\ref{mf:cor:luto} --- $H_N\le R\le N$: luto e exposição &
\pt{1}--\pt{3} & Lema~\ref{mf:lem:sub} (aqui) \\
Teor.~\ref{mf:thm:duas} --- a cisão é única: mede\,/\,ordena &
hipótese local & \code{fis:thm:duastorres}, \code{fis:thm:desce} \\
Cor.~\ref{mf:cor:ansiedade} --- a ansiedade vital não pára &
hipótese local & \code{fis:prop:cego} \\
Cor.~\ref{mf:cor:ponte} --- a passagem não custa &
--- & \code{fis:thm:ponte}, \code{fis:thm:H} \\
Teor.~\ref{mf:thm:outro} --- somar identifica; separar custa &
\pt{6} & \code{fis:thm:interfere}, \code{fis:thm:palavra} \\
Teor.~\ref{mf:thm:custo} --- apagar paga $\kappa\Phi$ &
\pt{1} + Landauer & \code{fis:thm:landauer}, \code{fis:cor:custofolga} \\
Teor.~\ref{mf:thm:rastro} --- o antiespaço é $\delta$, e halva &
--- & \code{fis:rastro}, \code{fis:thm:corte} \\
\end{longtable}}

\noindent\textbf{E a leitura desta tabela.} Dezoito enunciados, seis postulados,
uma hipótese local declarada e uma hipótese física importada e marcada.
\emph{Nenhuma outra coisa entrou.} O que a coluna do meio mostra é onde cada
teorema é frágil: quem negar \pt{4} perde a §\ref{mf:ancora} inteira e a
Rachadura~V; quem negar \pt{3} perde a §\ref{mf:adapta} e com ela o luto crónico.
\emph{É assim que se lê uma construção: pela coluna do meio.}

% ═══════════════════════════════════════════════════════════════════════════
% ═══════════════════════════════════════════════════════════════════════════
\part{As rachaduras}

\section{O dicionário, e o modo de ler a Parte II}\label{mf:dicionario}

\noindent Cada secção desta parte tem quatro blocos: \textbf{no livro} (as duas
citações, à letra), \textbf{porque não fecham}, \textbf{o fecho} (um enunciado da
Parte~I, pelo número) e \textbf{o compromisso}. A linha final resume os três.

\medskip
{\small
\begin{longtable}{@{}p{0.30\textwidth}p{0.40\textwidth}l@{}}
\toprule
no livro & é, nesta construção & onde \\
\midrule
\endfirsthead
\toprule
no livro & é, nesta construção & onde \\
\midrule
\endhead
\bottomrule
\multicolumn{3}{@{}r@{}}{\footnotesize\itshape continua na página seguinte}\\
\endfoot
\bottomrule
\endlastfoot
a experiência psíquica vivida & o índice $i\in I$, e a lista é ordenada & \pt{1} \\
a RPS na memória & a célula $x\in X$, com igualdade decidível e nada mais & \pt{1} \\
o fenómeno RAM & a realização $\pi\colon I\to X$: escreve, e não pergunta & \pt{1} \\
a disponibilidade histórica & $G(x)=\abs{\pi^{-1}(x)}$, a multiplicidade & \pt{2} \\
\cury{o passado não é lembrado, mas reconstruído} & a fibra perde-se, a medida conserva-se & Teor.~\ref{mf:thm:reconstroi} \\
o débito de toda recordação & a folga $\Phi$ --- \emph{medida} & Cor.~\ref{mf:cor:debito} \\
a âncora da memória & a \textbf{bola} de raio $r$ na árvore & \pt{4}, Teor.~\ref{mf:thm:ancora} \\
o deslocamento da âncora & mudar $r$, ou mudar de prefixo & Teor.~\ref{mf:thm:ancora}(4) \\
a autochecagem (o gatilho) & o incremento $a_x(t)$, que é $\mu G$ & Teor.~\ref{mf:thm:zetamu}(1) \\
o fenómeno do autofluxo & a metade que não mede & Cor.~\ref{mf:cor:ansiedade} \\
o \emph{eu} & a acumulação $\zeta$ --- memória, e não passo & \pt{5}, Teor.~\ref{mf:thm:zetamu}(2) \\
o gerenciamento do \emph{eu} & a deconvolução $\mu$ & Teor.~\ref{mf:thm:zetamu}(3) \\
o pensamento essencial & o campo: real, e no espaço & Teor.~\ref{mf:thm:duas} \\
o pensamento dialético & a metade \textbf{simétrica}: mede, desce & Teor.~\ref{mf:thm:duas}(1) \\
o pensamento antidialético & a metade \textbf{antissimétrica} & Teor.~\ref{mf:thm:duas}(2) \\
a leitura virtual (o «fenómeno LVD») & a transformada, e ela não paga & Cor.~\ref{mf:cor:ponte} \\
a psicoadaptação & $1/G$ & Teor.~\ref{mf:thm:adapta} \\
a ansiedade vital & a metade sem energia a descer & Cor.~\ref{mf:cor:ansiedade} \\
o caos psicodinâmico & apagar --- e apagar é o que custa & Teor.~\ref{mf:thm:custo} \\
o \cury{antiespaço} da consciência & $\delta$, medido, e halva & Teor.~\ref{mf:thm:rastro} \\
a traição inevitável da interpretação & a interferência: $G>1$ na soma & Teor.~\ref{mf:thm:outro} \\
\cury{colocar-se no lugar do outro} & o ramo: interferência por espaço & Teor.~\ref{mf:thm:outro}(3) \\
\cury{há apenas uma personalidade} & uma partição --- as bolas não se tocam & Teor.~\ref{mf:thm:ancora}(2) \\
\end{longtable}}

% ───────────────────────────────────────────────────────────────────────────
\section{Rachadura I --- a objecção à matemática destrói a própria teoria}\label{mf:r1}

\noindent\textbf{No livro} (cap.~\emph{A Metodologia\dots}, «As teorias na
Matemática também são limitadas»): \cury{até nas indiscutíveis operações de soma
há limitações, pois 1 mais 1 só é 2 se o primeiro 1 é, em todos os níveis
microessenciais, exatamente igual ao segundo 1. Porém, como tudo no universo é
essencialmente distinto\dots conclui-se que a Matemática\dots só é realmente
imutável por artifícios intelectuais que não se verificam no universo
essencial.}

\noindent E, no mesmo livro (cap.~\emph{A Memória e os Três Tipos de
Pensamentos}): \cury{cada informação ou experiência psíquica é armazenada na
memória como um sistema de código físico-químico, que chamo de representação
psicossemântica, RPS} --- e as RPSs \cury{são e serão sempre pobres, restritivas,
em relação à experiência original}.

\medskip\noindent\textbf{Porque não fecham.} A segunda frase \textbf{é} a
primeira operação que a primeira declarou impossível. Dizer que a RPS representa
a experiência de forma redutora é dizer que \emph{experiências distintas caem na
mesma representação} --- que é exactamente o que «$1+1=2$» afirma: dois objectos
distintos, uma contagem. Se a distinção microessencial proibisse a igualdade, não
haveria RPS nenhuma, porque não haveria \emph{a mesma} RPS a ser relida; e sem
isso não há memória, não há disponibilidade histórica, e não há teoria. \emph{O
livro destrói o instrumento de que precisa três capítulos adiante.}

\medskip\noindent\textbf{O fecho.} A objecção confunde \textbf{identidade dos
objectos} com \textbf{igualdade no contradomínio}, e \pt{1} só pede a segunda:
\[
\boxed{\;\begin{array}{c}
\pi\colon I\longrightarrow X\ \text{pede \textbf{só} igualdade decidível em }X\\[2pt]
\text{--- nem métrica, nem dimensão, nem estrutura ---}
\end{array}\;}
\]
Os dois uns \textbf{são} índices distintos, $i\neq j$; o que a igualdade afirma é
$\pi(i)=\pi(j)$, e daí não se conclui que $i$ e $j$ sejam o mesmo --- conclui-se
que estão na mesma fibra. E o Cor.~\ref{mf:cor:debito} dá o resto: \textbf{o que
Cury diz que a matemática esconde é precisamente o que $\Phi$ conta}. A frase
certa não é \emph{«$1+1=2$ é um artifício»}: é \emph{«$1+1=2$ afirma que a fibra
tem dois elementos»}, e a perda que ele quer denunciar tem número.

\medskip\noindent\textbf{O compromisso.} \pt{1}, e é barato: é menos do que o
livro já assume quando diz que a mesma RPS é relida.

\fecha{a igualdade é um artifício que o universo não verifica}
      {sem igualdade não há RPS, e a teoria da memória cai com ela}
      {\pt{1} pede só igualdade decidível, e a perda que ela produz é
       $\Phi$ --- Cor.~\ref{mf:cor:debito}}

% ───────────────────────────────────────────────────────────────────────────
\section{Rachadura II --- o \emph{eu} nasce na quinta etapa e opera na segunda}\label{mf:r2}

\noindent\textbf{No livro} (cap.~\emph{A Metodologia\dots}): as matrizes dos
pensamentos essenciais históricos \cury{são produzidas, em milésimos de segundos,
por quatro fenômenos que utilizam a história intrapsíquica arquivada na memória:
o fenômeno da autochecagem da memória, a âncora da memória, o fenômeno do
autofluxo e o eu} --- e a produção dessas matrizes é a \textbf{segunda} etapa da
interpretação.

\noindent E (cap.~\emph{O Gerenciamento do Eu\dots}): \cury{o eu é formado na
quinta etapa da interpretação, quando já foi desencadeada toda uma rica
construção de pensamentos nas etapas anteriores}; \cury{só quando o eu é gerado
na quinta etapa é que ele, retroativamente, pode exercer um gerenciamento da
construção de pensamentos}.

\medskip\noindent\textbf{Porque não fecham.} O \emph{eu} é ao mesmo tempo
\textbf{entrada da segunda etapa} e \textbf{saída da quinta}, no mesmo ciclo e no
mesmo instante. Não é realimentação bem posta --- para o ser, teria de haver dois
relógios, e o livro só tem um. E \cury{retroativamente} não é resposta: ou age
sobre uma segunda etapa \emph{futura}, e então não é retroactivo, é apenas tarde;
ou age sobre a que já passou, e então age sobre o que já não existe. \emph{O
livro usa a palavra e não escolhe o sentido.}

\medskip\noindent\textbf{O fecho.} O erro é de \textbf{nível}, e o
Teor.~\ref{mf:thm:zetamu} separa os dois:
\[
\boxed{\;G_t(x)=(a_x*\zeta)(t)\quad\text{(a escrita)}
\qquad\qquad
a_x(t)=G_t(x)-G_{t-1}(x)\quad\text{(a recuperação)}\;}
\]
Os três mordomos são o \textbf{incremento}; o \emph{eu} é a \textbf{acumulação};
e o gerenciamento não é faculdade nova --- é $\mu$. O
Cor.~\ref{mf:cor:vitima} dá a consequência que o livro observa e não deriva:
\cury{é mais fácil ele ser vítima do que agente modificador da sua história} é
$\zeta$ ser causal, e o atraso é \textbf{exactamente um passo} --- não é um
defeito de carácter. E \cury{retroativamente} ganha o sentido que faltava: não é
agir no passado, é \textbf{ler o passado do campo presente}.

\medskip\noindent\textbf{O compromisso.} \pt{5} tira ao \emph{eu} o estatuto de
\emph{fenómeno que lê}, que o livro lhe dá quatro vezes. Em troca dá-lhe uma
operação exacta e um lugar sem circularidade. \emph{Não se pode ter as duas
coisas.}

\fecha{o \emph{eu} é um dos quatro leitores, e forma-se depois dos outros três}
      {entrada da segunda etapa e saída da quinta, no mesmo instante}
      {Teor.~\ref{mf:thm:zetamu}: os três são $\mu$, o \emph{eu} é $\zeta$, e
       gerir é desacumular; Cor.~\ref{mf:cor:vitima} dá o atraso de um passo}

% ───────────────────────────────────────────────────────────────────────────
\section{Rachadura III --- quatro leitores, e nenhuma assinatura}\label{mf:r3}

\noindent\textbf{No livro} (cap.~\emph{A Memória e os Três Tipos de
Pensamentos}): \cury{quando o eu dá continuidade às cadeias psicodinâmicas
produzidas pelos demais fenômenos que lêem a memória, ele o faz através de uma
construção mais analítica e histórico-crítica} --- o que exige que o \emph{eu}
saiba \textbf{qual} dos quatro produziu a cadeia que tem à frente.

\noindent E (cap.~\emph{O Fenômeno da Autochecagem}): o registo é feito
\cury{por um fenômeno que chamo de fenômeno RAM --- registro automático da
memória}, e é \cury{involuntário e automático}. O RAM regista o pensamento;
\emph{não regista quem o construiu}.

\medskip\noindent\textbf{Porque não fecham.} O livro pede ao \emph{eu} uma
distinção que os seus próprios postulados tornam impossível. Se as quatro
leituras produzem matrizes essenciais e o RAM as regista todas do mesmo modo,
então \textbf{nada na memória diz de onde veio o quê}. \emph{E o livro precisa
dessa separação em todos os capítulos clínicos.}

\medskip\noindent\textbf{O fecho, e é o mais duro deste documento: a distinção é
provadamente impossível, e a que é possível está identificada.} O
Teor.~\ref{mf:thm:multifocal} diz que $G$ é invariante sob qualquer reatribuição
de agentes, e que \textbf{nenhuma função do campo} distingue autores; o
Teor.~\ref{mf:thm:folha} diz o que se repõe e o que não se repõe:

\begin{center}\small
\begin{tabular}{@{}lll@{}}
\toprule
o que se quer separar & a coordenada que o faz & sai do campo? \\
\midrule
\textbf{a ordem} da visita & a folha $k(i)$ & \textbf{sim} \\
\textbf{o agente} que escreveu & o identificador & \textbf{não} \\
\bottomrule
\end{tabular}
\end{center}

\noindent E o mais importante: \textbf{isto não enfraquece a tese do livro ---
substitui-a por uma melhor}. A inteligência não é multifocal porque há quatro
fenómenos; é multifocal porque o campo não tem coluna para o autor. \emph{A
contagem quatro é da observação; o teorema não precisa dela.}

\medskip\noindent\textbf{O compromisso.} Perde-se a frase segundo a qual o
\emph{eu} trata criticamente «as cadeias dos outros fenómenos». Ganha-se uma
operação que se enuncia e se mede, e uma explicação para uma coisa que o livro
observa sem explicar: porque é que a introspecção honesta produz
\emph{cronologia} e nunca produz \emph{atribuição}.

\fecha{o \emph{eu} distingue as cadeias que os outros três produziram}
      {nada no registo carrega o autor: o RAM regista o pensamento, não o
       construtor}
      {Teor.~\ref{mf:thm:multifocal}: nenhuma função do campo distingue autores;
       Teor.~\ref{mf:thm:folha}: repõe-se a ordem, não a autoria}

% ───────────────────────────────────────────────────────────────────────────
\section{Rachadura IV --- não há lembrança, e mesmo assim há relação lógica}\label{mf:r4}

\noindent\textbf{No livro} (cap.~\emph{A Memória e os Três Tipos de
Pensamentos}): \cury{não existe lembrança. Não há lembrança das informações
contidas na memória, mas reconstrução das mesmas}; \cury{a história existencial
(intrapsíquica) está morta essencialmente na memória}.

\noindent E (cap.~\emph{O Fenômeno da Autochecagem}): \cury{embora a construção
de pensamentos ultrapasse os limites da lógica, ela possui um sistema de relação
lógica; caso contrário, não seria possível produzir ciência, não haveria verdade
científica}.

\medskip\noindent\textbf{Porque não fecham.} Perda total e fidelidade suficiente
são afirmadas lado a lado, sem \textbf{nenhuma grandeza} que as separe. O livro
diz \cury{toda «recordação» tem um débito emocional em relação à experiência
original} --- e a palavra \emph{débito} promete um número que nunca aparece.

\medskip\noindent\textbf{O fecho, e tem duas partes.} \emph{Primeiro}, o
Teor.~\ref{mf:thm:reconstroi} separa o que se conserva do que se perde:
\[
\boxed{\;\sum_x G(x)=\abs{I}\ \text{sempre}
\qquad\text{e}\qquad
\text{a medida \textbf{não distingue} realizações}\;}
\]
\emph{É a frase que o livro procurava}: a memória conserva a \textbf{contagem} e
perde a \textbf{fibra}. \emph{Segundo}, o Cor.~\ref{mf:cor:debito} dá o número:
$\Phi=\abs{I}-\abs{X}$. E o Teor.~\ref{mf:thm:rastro} dá a forma da recordação:
o par $(m,\delta)$ --- o que se guarda, e o que se perdeu, medido. \emph{O
«débito emocional» de Cury é $\delta$.} A sua observação de que a mesma perda é
recordada com intensidades diferentes é o Cor.~\ref{mf:cor:folhas}: \emph{não é a
mesma recordação duas vezes; são duas folhas da mesma célula}.

\medskip\noindent\textbf{O compromisso.} Aceita-se que o que a memória guarda é
um par e não um conteúdo --- o que o livro, de resto, já negava --- e paga-se com
a única coisa que ele queria e não tinha: \emph{a medida do que se perdeu}.

\fecha{a memória não devolve, reconstrói; e mesmo assim há relação lógica}
      {perda total e fidelidade suficiente, sem grandeza que as separe}
      {Teor.~\ref{mf:thm:reconstroi} e Cor.~\ref{mf:cor:debito}: conserva-se a
       medida, perde-se a fibra, e o débito é $\Phi$}

% ───────────────────────────────────────────────────────────────────────────
\section{Rachadura V --- a âncora, que o autor declara não saber definir}\label{mf:r5}

\noindent\textbf{No livro} (cap.~\emph{A Âncora da Memória}), e a citação é a
confissão: \cury{a âncora da memória é difícil de ser definida, mas em síntese
ela se refere a um foco ou «território» de leitura da memória num determinado
momento da existência}.

\noindent E, no mesmo capítulo, a âncora tem de fazer \textbf{três} coisas
precisas: delimitar o que está disponível; ser \textbf{uma} em cada momento (é
ela que faz do mesmo homem um \cury{anjo social} e um \cury{carrasco da
família}); e poder ser \textbf{alargada} pelo \emph{eu}.

\medskip\noindent\textbf{Porque não fecha.} Um \cury{território} sem critério não
delimita nada. O livro não diz o que faz de um conjunto de RPSs um território,
nem porque é que os territórios não se sobrepõem, nem porque é que há exactamente
um de cada vez --- e as três coisas são carga estrutural: sem elas, a explicação
do pânico, da embriaguez e do \cury{cárcere intelectual} não tem sobre o que
assentar. \emph{Uma peça que suporta o edifício não pode entrar declarada como
indefinível.}

\medskip\noindent\textbf{O fecho.} \pt{4} define-a, e o
Teor.~\ref{mf:thm:ancora} demonstra as três exigências, uma por cláusula:
delimita (é uma célula), é uma (as bolas particionam, e daí o
Cor.~\ref{mf:cor:umsotem}), e alarga-se de dentro (todo ponto é centro). E dá uma
quarta coisa que o livro não tinha: \textbf{são dois deslocamentos e não um} ---
mudar o raio, ou mudar de prefixo. É por isso que o pânico \emph{estreita} e a
festa de formatura \emph{muda de sítio}: \emph{são movimentos diferentes}. O
Cor.~\ref{mf:cor:prox} acrescenta o número: a proximidade entre territórios é
$(2q-N)/N$, crescente no prefixo partilhado.

\medskip\noindent\textbf{O compromisso.} \pt{4}, e é o mais caro do documento.
Em troca, as três propriedades que o livro precisava de postular tornam-se
cláusulas de um teorema --- e a euclidiana, a alternativa óbvia, \emph{falha
ambas}, o que está medido.

\fecha{a âncora é um «território» de leitura, e o autor declara não a saber
       definir}
      {sem critério, não delimita, não explica a unicidade e não explica o
       alargamento}
      {\pt{4} + Teor.~\ref{mf:thm:ancora}: é a bola, as bolas particionam, todo
       ponto é centro, e há dois deslocamentos}

% ───────────────────────────────────────────────────────────────────────────
\section{Rachadura VI --- o RAM reforça o que a psicoadaptação apaga}\label{mf:r6}

\noindent\textbf{No livro} (cap.~\emph{A Reorganização do Caos\dots}):
\cury{quanto mais emoção tensa, ansiosa ou prazerosa as cadeias de pensamentos
possuem, mais o fenômeno RAM as registra em zonas privilegiadas da memória, o que
facilita sua leitura}.

\noindent E (cap.~\emph{A Comunicação Social Mediada\dots}): a psicoadaptação é
\cury{uma «incapacidade» da energia emocional e motivacional de se submeter aos
mesmos processos de transformações essenciais, diante da exposição das matrizes
idênticas ou semelhantes de pensamentos essenciais históricos}; \cury{todos os
estímulos que interpretamos, sejam dolorosos ou prazerosos, com o passar do tempo
não provocam a mesma emoção}.

\medskip\noindent\textbf{Porque não fecham.} A \emph{mesma} repetição faz as duas
coisas contrárias: torna a experiência \textbf{mais disponível} e \textbf{menos
sentida}. O livro afirma as duas sem lei conjunta, e usa-as em sentidos opostos
no mesmo parágrafo --- o luto crónico é explicado por a multiplicidade das
matrizes \emph{derrotar} a psicoadaptação, o que só faz sentido se houver uma
conta, e a conta não está lá.

\medskip\noindent\textbf{O fecho, e é o melhor do documento.} O
Teor.~\ref{mf:thm:adapta} mostra que são duas leituras de um campo, e que a
segunda é uma divisão:
\[
\boxed{\;
\underbrace{G\ \text{sobe}}_{\text{o RAM: acumula}}
\qquad
\underbrace{\text{o incremento vale }1}_{\text{sempre}}
\qquad
\underbrace{\frac{G+1}{G}-1=\frac1G}_{\text{o que se sente}}
\;}
\]
\emph{Não há contradição nenhuma}: uma frase é sobre o campo, a outra sobre o seu
gradiente, e a segunda é o inverso da primeira. E o
Cor.~\ref{mf:cor:luto} dá o mecanismo do luto crónico, que o livro observa e não
deriva: com $N$ visitas, a resposta total obedece a $H_N\le R\le N$, e
\textbf{repartir maximiza, concentrar minimiza}. Não se sofre mais: sofre-se
\emph{em mais lados}. E a recíproca é a exposição.

\medskip\noindent\textbf{O compromisso.} \pt{3} --- identificar «o que se sente»
com o gradiente e não com o campo. \emph{Quem preferir identificá-lo com $G$ tem
de explicar porque é que o vigésimo dia com o carro novo não vale vinte vezes o
primeiro.}

\fecha{a repetição regista mais fundo, e a repetição faz sentir menos}
      {a mesma causa com dois efeitos contrários, e sem lei que os ligue}
      {Teor.~\ref{mf:thm:adapta}: o incremento vale $1$ e o que se sente é
       $1/G$; Cor.~\ref{mf:cor:luto}: $H_N\le R\le N$}

% ───────────────────────────────────────────────────────────────────────────
\section{Rachadura VII --- o fluxo que nunca repousa, e a terapia que o faz repousar}\label{mf:r7}

\noindent\textbf{No livro} (cap.~\emph{O Fenômeno do Autofluxo\dots}): \cury{o
campo de energia psíquica nunca se «equilibra», mas vive continuamente um
«desequilíbrio psicodinâmico»}; e os termos correntes da psicologia ---
\cury{equilíbrio emocional}, \cury{pessoa equilibrada} --- \cury{são
superficiais; não procedem do ponto de vista científico}.

\noindent E (cap.~\emph{O Gerenciamento do Eu\dots}): \cury{desacelerar o fluxo
dos pensamentos e alargar os territórios de leitura da memória (âncora) é
fundamental para que o eu possa administrar a produção dos pensamentos com
lógica}.

\medskip\noindent\textbf{Porque não fecham.} Se o campo \textbf{não pode}
repousar --- por princípio, e o livro chama-lhe \emph{princípio dos princípios}
--- então desacelerar até uma leitura organizada é impossível, e a terapia que o
livro pratica não tem sobre o que operar. Se pode, o \cury{nunca} é falso e os
termos que ele recusa eram os certos.

\medskip\noindent\textbf{O fecho.} O Teor.~\ref{mf:thm:duas}: a cisão
$W=W_s+W_a$ é \textbf{única}, e as duas metades fazem coisas disjuntas --- uma
mede e pára em passos finitos, a outra não tem forma quadrática e portanto não
tem função de Lyapunov. \emph{As duas afirmações de Cury são ambas verdadeiras,
de metades diferentes}: a ansiedade vital é a antissimétrica
(Cor.~\ref{mf:cor:ansiedade}), o repouso da terapia é a simétrica.

\medskip\noindent\textbf{E o que o livro afirma a mais, e que não se segue.} Cury
faz do desequilíbrio um princípio \emph{produtivo}. A cláusula~(3) do
Teor.~\ref{mf:thm:duas} não dá isso: a antissimetria \textbf{tira} a função de
Lyapunov e \textbf{não põe} um ciclo no lugar dela. \emph{E a obra pagou por o
ter aprendido:} o medidor que afirmava o contrário falha, e a falha é o resultado.

\medskip\noindent\textbf{O compromisso.} A hipótese local do
Teor.~\ref{mf:thm:duas} --- que a leitura se represente por uma matriz com
transposta. É real, e é a mais forte da Parte~I depois de \pt{4}.

\fecha{o campo nunca repousa, por princípio}
      {e a terapia do autor consiste em o fazer repousar}
      {Teor.~\ref{mf:thm:duas}: a cisão é única --- uma metade mede e pára, a
       outra ordena e não tem energia a descer}

% ───────────────────────────────────────────────────────────────────────────
\section{Rachadura VIII --- três pensamentos, e um fenómeno tapa-buraco}\label{mf:r8}

\noindent\textbf{No livro} (cap.~\emph{A Memória e os Três Tipos de
Pensamentos}): há três tipos, e só três --- o essencial (inconsciente, e \cury{de
natureza real}), o dialético e o antidialético.

\noindent E, para explicar a passagem do primeiro aos outros dois, entra um
quarto objecto: \cury{o processo de leitura virtual das matrizes dos pensamentos
essenciais é realizado por um fenômeno intrapsíquico, chamado de fenômeno LVD}. E
logo a seguir, para o antidialético: \cury{os pensamentos antidialéticos são
produzidos por uma leitura virtual mais difusa, provavelmente por outro fenômeno
que não o LVD, ou então por um processo de leitura difusa desse fenômeno}.

\medskip\noindent\textbf{Porque não fecha.} A contagem «três» não é derivada de
nada: para a sustentar entra um fenómeno com iniciais, e para sustentar
\emph{esse} entra um \cury{provavelmente} entre duas hipóteses que o autor não
escolhe. \emph{Uma teoria que precisa de um fenómeno novo para cada passagem não
tem três tipos: tem tantos quantos os buracos.}

\medskip\noindent\textbf{O fecho, e ele desmonta o quarto objecto em vez de o
substituir.} O Teor.~\ref{mf:thm:duas} dá \textbf{um} campo e \textbf{duas}
metades, com a cisão única; e a caracterização do antidialético cai no sítio pelo
Cor.~\ref{mf:cor:ansiedade} --- \emph{o que não tem energia não tem como ser
dito, e continua a mandar}, que é à letra o que Cury descreve. E o «LVD» não é
substituído: é \textbf{dispensado} pelo Cor.~\ref{mf:cor:ponte}, porque a
travessia não paga. \emph{Um fenómeno é preciso onde há custo por justificar;
aqui o custo é zero, e o \cury{provavelmente} desaparece com a pergunta que o
gerou.}

\medskip\noindent\textbf{O compromisso.} Perde-se a arrumação em «três tipos».
Ganha-se a dispensa de dois fenómenos postulados e de um \cury{provavelmente}.

\fecha{três tipos, ligados por um «fenómeno LVD» e por um segundo fenómeno
       provável}
      {a contagem não é derivada, e cada passagem exige um fenómeno novo}
      {Teor.~\ref{mf:thm:duas}: um campo e duas metades; Cor.~\ref{mf:cor:ponte}:
       a passagem não paga, logo não precisa de fenómeno}

% ───────────────────────────────────────────────────────────────────────────
\section{Rachadura IX --- uma distância infinita que mesmo assim se gradua}\label{mf:r9}

\noindent\textbf{No livro} (cap.~\emph{A Metodologia\dots}): \cury{entre o
conhecimento consciente e a realidade essencial dos objetos e fenômenos há uma
distância infinita, um antiespaço entre o virtual (consciência) e o real
(fenômeno essencial)}; a consciência \cury{nunca incorpora a realidade essencial
dos fenômenos e dos objetos}.

\noindent E, poucas páginas adiante: \cury{quanto mais a consciência existencial
mantém um sistema de relação com a verdade essencial, mais ela conquista um
status de verdade científica, capaz de resultar em aplicabilidade e
previsibilidade de fenômenos}.

\medskip\noindent\textbf{Porque não fecham.} Uma distância infinita não admite
\emph{mais} nem \emph{menos}. Se o antiespaço é intransponível e sempre o mesmo,
não há como uma consciência estar mais perto da verdade essencial do que outra --
e cai com isso a única defesa que o livro tem contra o relativismo que ele
próprio recusa em voz alta. \emph{O autor precisa de uma medida e declarou-a
infinita.}

\medskip\noindent\textbf{O fecho.} O Teor.~\ref{mf:thm:rastro}: o que fica é o
par $(m,\delta)$; a consciência guarda $m$ e \textbf{não} incorpora o ponto ---
verdade; e o que ficou de fora é $\delta$, que é um número e que desce com
largura seguinte $\le\delta^2/m$ e cota logarítmica. \emph{O que é infinito é o
número de batidas, e não a distância.} E há um desfecho para a tese forte do
livro: pela cláusula~(3), o ponto \textbf{pode não caber em grade nenhuma} --- e
mesmo então as duas faces cercam-no, com medida.

\medskip\noindent\textbf{E o método do livro sai justificado por um argumento que
ele não tinha}: o Cor.~\ref{mf:cor:triade} mostra que o encaixe exige bater
\textbf{alternadamente} nas duas faces. \emph{A tríade pergunta\,/\,dúvida\,/\,
crítica é a alternância}, e é essa a sua razão estrutural.

\medskip\noindent\textbf{O compromisso.} Troca-se um absoluto retórico por uma
grandeza. Quem quiser guardar o absoluto guarda-o --- mas tem de riscar a frase
sobre a verdade científica, e com ela a distinção entre uma teoria e um palpite.

\fecha{entre a consciência e a coisa há uma distância infinita}
      {e mesmo assim umas consciências estão mais perto da verdade do que
       outras}
      {Teor.~\ref{mf:thm:rastro}: a distância é $\delta$, mede-se e halva; o
       infinito é o número de batidas}

% ───────────────────────────────────────────────────────────────────────────
\section{Rachadura X --- a traição inevitável da interpretação, sem mecanismo}\label{mf:r10}

\noindent\textbf{No livro} (cap.~\emph{A Comunicação Social Mediada\dots}):
\cury{não comunicamos a essência do que pensamos e sentimos}; o observador
\cury{reconstrói interpretativamente} a experiência do outro, e \cury{conhecemos
as pessoas a partir de nós mesmos}. A isto ele chama a \cury{traição inevitável
da interpretação}.

\noindent E, no mesmo capítulo, o mecanismo: a autochecagem lê \textbf{ao mesmo
tempo} as RPS \emph{diretivas} e as \emph{associativas}, e das duas sai
\emph{uma} matriz de pensamentos essenciais históricos.

\medskip\noindent\textbf{Porque não fecha.} \cury{Inevitável} é uma palavra forte
e o livro não a sustenta: nada diz \emph{porque} é que a reconstrução tem de
distorcer, nem \emph{quanto}, nem \emph{o que fazer para distorcer menos} ---
embora exija que se faça. A prescrição (\cury{aprender a se colocar no lugar do
outro}) fica sem custo declarado, e por isso sem força: \emph{uma recomendação
que não custa nada não explica porque é que quase ninguém a cumpre}.

\medskip\noindent\textbf{O fecho.} O Teor.~\ref{mf:thm:outro}: a leitura do outro
é uma \textbf{soma}, a soma tem \textbf{núcleo}, e dois «outros» distintos podem
produzir a mesma leitura --- \emph{a interferência é $G>1$}. É esse o mecanismo
que faltava, e é consequência de somar, não fraqueza moral. O
Cor.~\ref{mf:cor:outroeu} dá o \emph{quanto}: o outro entra como incremento sobre
um campo que já carrega a história do observador, e o que dele se lê é $1/G$ ---
\emph{quanto mais história, menos o outro aparece}. E a prescrição fica com o
preço: a cláusula~(3) desfaz a interferência com \textbf{uma} coordenada --- o
ramo --- e a cláusula~(4) diz que ela \textbf{ocupa lugar}.

\medskip\noindent\textbf{O compromisso.} \pt{6}, e com ele uma consequência
incómoda: \emph{a distorção não se elimina com boa vontade}; elimina-se com
espaço, ou não se elimina.

\fecha{a interpretação do outro trai sempre, e deve-se aprender a colocar-se no
       lugar dele}
      {nenhum mecanismo diz porquê, e a prescrição não tem custo declarado}
      {Teor.~\ref{mf:thm:outro}: somar identifica e separar custa espaço;
       Cor.~\ref{mf:cor:outroeu} dá o quanto, e é $1/G$}

% ───────────────────────────────────────────────────────────────────────────
\section{Rachadura XI --- o caos é o motor, e ninguém paga a conta}\label{mf:r11}

\noindent\textbf{No livro} (cap.~\emph{A Memória e os Três Tipos de
Pensamentos}): \cury{a morte do presente é a única possibilidade de expansão do
próprio presente}; \cury{tem de haver a morte ou o caos do «eu sou»\dots para que
surja o «fui histórico»}.

\noindent E (cap.~\emph{A Reorganização do Caos\dots}): esse ciclo é
\cury{extremamente rápido}, corre \cury{milhares de vezes ao dia}, e é
apresentado como a condição de toda a construção --- \emph{sem uma linha sobre o
que custa}.

\medskip\noindent\textbf{Porque não fecha.} O livro faz do apagamento a fonte de
tudo e trata-o como gratuito. Mas também afirma que o campo psíquico
\cury{co-interfere} com o campo físico-químico do cérebro e se transmuta nele ---
e no momento em que uma coisa se transmuta em física, apagar deixa de ser de
graça. \emph{A teoria pede uma economia e não abre livro nenhum.}

\medskip\noindent\textbf{O fecho.} O Teor.~\ref{mf:thm:custo}: o piso do
invertível é \textbf{zero exacto}, o destrutivo paga por cada bit, e a conta
total é $\kappa\Phi$ --- \emph{a mesma folga da Rachadura~IV, vezes uma
constante}. E a cláusula~(4) dá a frase que reordena o capítulo:
\[
\boxed{\;\text{o dual guardado é o piso não pago}\;}
\]
\textbf{Donde a prescrição central de Cury fica certa, e pela razão que ele não
deu.} Ele defende que se \cury{revise, recicle, reorganize} em vez de se
descartar, e apresenta-o como higiene intelectual. Aqui é uma factura: \emph{o
que se guarda reverte a custo zero; o que se deita fora é cobrado.}

\medskip\noindent\textbf{O compromisso.} Importa-se uma hipótese física, e ela
vai marcada. E aceita-se que o \cury{caos psicodinâmico} é apagamento --- que é a
leitura do livro, mas é uma leitura: quem preferir chamar-lhe
\emph{recodificação} não paga nada e também não explica nada.

\fecha{o caos do presente é a condição de toda a construção}
      {apagar é apresentado como gratuito, num campo que se transmuta em física}
      {Teor.~\ref{mf:thm:custo}: piso zero exacto para o invertível,
       $\kappa\Phi$ para o destrutivo, e o dual guardado é o piso não pago}

% ───────────────────────────────────────────────────────────────────────────
\section{Rachadura XII --- escravo da âncora, e capaz de a alargar}\label{mf:r12}

\noindent\textbf{No livro} (cap.~\emph{A Âncora da Memória}): \cury{muitas vezes
achamos que somos livres na produção das idéias, mas somos prisioneiros da âncora
da memória}; \cury{quem determina o cárcere intelectual é a âncora da memória}.

\noindent E, três parágrafos depois: \cury{devemos aprender a gerenciá-la,
reciclá-la criticamente e deslocá-la com liberdade}; e nos focos de tensão,
\cury{se ele não tiver êxito em abrir a âncora e alargar os territórios de
leitura da memória a produção de pensamentos e reações emocionais pode ser
insegura, fóbica, agressiva}.

\medskip\noindent\textbf{Porque não fecham.} Se tudo o que se lê vem de dentro do
território ancorado, \textbf{a decisão de alargar também tem de vir de lá} --- e
o que está lá é, por hipótese, o que produz a reacção estreita. É um regresso:
para sair é preciso já estar fora. O livro nota o problema (\cury{o eu tanto
dirige os deslocamentos da âncora da memória como é dirigido por esses
deslocamentos}) e deixa-o como paradoxo.

\medskip\noindent\textbf{O fecho, e tem três peças.} \emph{Primeira}: pela
cláusula~(3) do Teor.~\ref{mf:thm:ancora}, \textbf{todo ponto é centro}, e a bola
de raio maior tem o mesmo ponto por centro --- \emph{subir $r$ é uma operação
disponível de dentro}, e o regresso desfaz-se. \emph{Segunda}: pela
cláusula~(1) do Teor.~\ref{mf:thm:duas} a descida é bem-fundada e \textbf{pára}
--- o \cury{stop introspectivo} acaba em número finito de passos. \emph{Terceira},
e é a que dá o mecanismo do \cury{resgate da liderança do eu nos focos de
tensão}: o passo natural vai ao \textbf{menor} $G$ ---
\fis{fis:thm:sinal}{Teor. «o passo foge do campo»} --- e deliberar é ir ao
$\arg\max$, isto é, \emph{voltar} ao sítio onde já se esteve muitas vezes.
\textbf{E isso não é hipótese nova}: o passo do $\arg\max$ é a imagem do passo do
$\arg\min$ pelo operador $\mathcal{D}$, sem parâmetro a escolher
(\fis{fis:thm:forca}{Teor. «a força é a dobra da queda»}). O custo vem no mesmo
teorema: atravessar uma célula já visitada custa
$\partial\tau=\sqrt{G(x)h_{ii}}$ --- \emph{voltar demora, e demora
$\sqrt G$}. É a razão estrutural de a introspecção ser lenta em proporção ao que
já lá está.

\medskip\noindent\textbf{O compromisso.} Fica-se com uma liberdade modesta: não é
escolher o que se pensa --- é escolher \textbf{o sentido do passo}, e nada mais.
É menos do que o livro promete nos capítulos finais, é mais do que ele sustenta
nos teóricos, e \emph{é a única que não faz regresso}.

\fecha{somos prisioneiros da âncora, e devemos aprender a alargá-la}
      {a decisão de alargar teria de ser lida de dentro do que a estreita ---
       regresso}
      {Teor.~\ref{mf:thm:ancora}(3): todo ponto é centro, logo alarga-se de
       dentro; e deliberar é $\arg\max$ contra $\arg\min$}

% ───────────────────────────────────────────────────────────────────────────
\section{Rachadura XIII --- uma só mente, e dois territórios que não se tocam}\label{mf:r13}

\noindent\textbf{No livro} (cap.~\emph{A Âncora da Memória}): há pessoas que fora
de casa são \cury{sorridentes, solícitas, pacientes e sóbrias (um «anjo
social»)} e dentro de casa \cury{se transformam em «carrascos da família»}. E o
veredicto: \cury{o termo dupla personalidade é dialeticamente pobre e
cientificamente inverificável. Há apenas uma personalidade porque há apenas uma
mente, apenas um campo de energia psíquica}.

\medskip\noindent\textbf{Porque não fecha.} O argumento é de \textbf{substância}
(há um só campo) contra uma observação de \textbf{estrutura} (há duas produções
que não comunicam). Não responde: um só campo pode ter partes que não se
alcançam, e é isso, e não a duplicação da substância, o que o fenómeno pede.
Pior: o próprio livro fornece o que faltava --- os territórios da âncora --- e
não o usa aqui, porque nunca disse que os territórios são \emph{disjuntos}.

\medskip\noindent\textbf{O fecho.} A cláusula~(2) do Teor.~\ref{mf:thm:ancora}:
duas bolas do mesmo raio \textbf{ou coincidem ou são disjuntas}, e para cada $r$
formam uma \textbf{partição} de $X$:
\[
\boxed{\;\text{um só } X\ \text{---}\ \text{e as suas bolas não se tocam}\;}
\]
\emph{Um campo, e uma partição dele.} Não é preciso duplicar mente nenhuma: é
preciso notar que uma partição de um conjunto não faz dois conjuntos. E a obra
tem o enunciado exacto para o par: o
\fis{fis:thm:duoesp}{Teor. «o par é duomorfo, e a interface é não se tocarem»}
--- \emph{a interface é a ausência de contacto}.

\medskip\noindent\textbf{Donde o veredicto de Cury fica certo, e o argumento
muda.} \cury{Não há dupla personalidade} --- verdade, e agora por uma razão.
\cury{Porque há apenas uma mente} --- irrelevante, e é isso que torna o argumento
fraco: a unicidade da substância não implicaria nada sobre a comunicação entre as
partes. \emph{O que faz o trabalho não é haver um campo; é as bolas
particionarem.}

\medskip\noindent\textbf{O compromisso.} Passa-se de uma tese sobre \emph{o que a
mente é} para uma tese sobre \emph{como ela está partida} --- e a segunda é
verificável e a primeira não.

\fecha{não há dupla personalidade porque há uma só mente}
      {argumento de substância contra uma observação de estrutura --- não
       responde}
      {Teor.~\ref{mf:thm:ancora}(2): um só $X$, e as suas bolas ou coincidem ou
       são disjuntas}

% ───────────────────────────────────────────────────────────────────────────
\section{Rachadura XIV --- «os computadores jamais», sem critério}\label{mf:r14}

\noindent\textbf{No livro} (cap.~\emph{Minha Trajetória de Pesquisa}): \cury{será
impossível que um dia os computadores conquistem essa consciência}; \cury{para os
computadores, o tudo e o nada, o ter e o ser, um segundo e a eternidade, serão
inevitavelmente sempre a mesma coisa}.

\noindent E o motivo, dado dois capítulos adiante: \cury{nos computadores
procuramos as informações através de rígidos e engessados sistemas de códigos},
ao passo que na memória humana \cury{ela não é lida por índice, mas por conteúdos
seletivos}.

\medskip\noindent\textbf{Porque não fecham.} O motivo é \textbf{arquitectural} e
a conclusão é \textbf{metafísica}. E o motivo é negado pelo próprio livro: as
RPSs são \cury{sistemas de códigos físico-químicos} arquivados em
\cury{determinados sítios do córtex cerebral} e lidos por endereço associativo
--- que é uma arquitectura de índice com outro nome. \emph{Se o critério é a
arquitectura, não separa o que o autor quer separar; e se não é, o autor não deu
critério nenhum.}

\medskip\noindent\textbf{O fecho: o critério existe, é estrutural, e não é sobre
substrato.}
\[
\boxed{\;\text{a memória está \textbf{no agente}}
\qquad\text{contra}\qquad
\text{a memória está \textbf{no espaço}}\;}
\]
É a cláusula~(3) do \fis{fis:thm:mult}{Teor. «multiplicidade geométrica»} ---
\emph{nenhum passo consulta $\pi(0),\dots,\pi(t)$} --- que a §\ref{mf:alicerce}
importou e que a cláusula~(1) do Teor.~\ref{mf:thm:zetamu} usa. E a diferença tem
consequência operacional, que se mede: pelo
\fis{fis:thm:periodo}{Teor. «o período lê-se no campo»}, com a memória no espaço
\textbf{um cursor basta} --- a cauda e o período lêem-se no campo, sem segunda
passagem e sem comparar posições, onde o algoritmo clássico precisa de dois.

\medskip\noindent\textbf{Donde a afirmação de Cury tem de ser reescrita, e fica
mais fraca e verdadeira.} O que a máquina da arquitectura clássica não tem não é
consciência: é um \textbf{campo que ela não carrega}. E no momento em que passa a
carregá-lo, a distinção deixa de ser sobre silício e passa a ser sobre onde está
a marca --- propriedade de \emph{desenho}, não de matéria. \emph{O «jamais» não
se segue de nada que o livro tenha estabelecido.}

\medskip\noindent\textbf{E note-se o que este documento \emph{não} está a
afirmar.} Não se afirma que uma máquina com o campo no espaço tenha consciência
existencial; afirma-se que o argumento do livro não estabelece o contrário, e que
o critério que ele procurava existe e é este. \emph{Trocar uma impossibilidade
por uma pergunta em aberto é um ganho, e não uma concessão.}

\fecha{os computadores jamais terão consciência existencial, porque lêem por
       índice}
      {motivo arquitectural para conclusão metafísica --- e o motivo é negado
       pelas próprias RPSs do livro}
      {memória no agente contra memória no espaço --- e mede-se: um cursor onde
       o clássico precisa de dois}

% ═══════════════════════════════════════════════════════════════════════════
\section{O que fica de pé}\label{mf:pe}

\noindent\textbf{E fica muito, e convém dizê-lo com a mesma clareza com que se
disseram as rachaduras.} Um livro com catorze junções partidas continua a poder
estar certo no essencial, e este está --- e está de um modo raro: as suas teses
centrais, que na altura eram heresias, saem daqui \textbf{demonstradas} em vez de
observadas.

\medskip
{\small
\begin{tabular}{@{}p{0.44\textwidth}p{0.49\textwidth}@{}}
\toprule
a tese do livro & e aqui é \\
\midrule
a construção da inteligência é \textbf{multifocal} --- não há autor único &
Teor.~\ref{mf:thm:multifocal}. \emph{E melhor do que ele a punha: não é que haja
quatro; é que o campo não tem coluna para eles} \\[4pt]
o \emph{eu} \textbf{não é} o líder dos processos que o produzem &
Cor.~\ref{mf:cor:vitima}: uma acumulação não pode ser um incremento, e o atraso
é de um passo \\[4pt]
o passado \textbf{não é lembrado}, é reconstruído &
Teor.~\ref{mf:thm:reconstroi} + Cor.~\ref{mf:cor:debito}: conserva-se a medida,
perde-se a fibra, e o que se perdeu tem número \\[4pt]
a leitura da memória é \textbf{inevitável} diante do estímulo &
a escrita é a realização, e ela não pergunta: \pt{1} com
Teor.~\ref{mf:thm:zetamu}(1) \\[4pt]
todos têm \textbf{o mesmo espectáculo} nos bastidores &
Teor.~\ref{mf:thm:multifocal}(2): nenhuma função do campo distingue autores ---
\emph{nem, portanto, pessoas} \\[4pt]
a dor perpetua-se por \textbf{multiplicidade} de experiências &
Cor.~\ref{mf:cor:luto}: $H_N\le R\le N$, e repartir maximiza \\[4pt]
\textbf{reciclar} vale mais do que descartar &
Teor.~\ref{mf:thm:custo}(4): o dual guardado é o piso não pago \\[4pt]
a tríade pergunta\,/\,dúvida\,/\,crítica, e a alternância &
Cor.~\ref{mf:cor:triade}: o encaixe exige as \emph{duas} faces \\
\bottomrule
\end{tabular}}

\medskip\noindent\textbf{E há uma que merece linha própria, porque é a mais
bonita.} Cury faz um argumento humanista a partir da construção dos pensamentos:
como todos leem a memória do mesmo modo, \cury{não deveria haver uma
intelligentsia como classe, ou qualquer outra forma de classificação intelectual
que distinga os seres humanos}. \emph{Aqui isso é um teorema, e não uma
exortação}: pela cláusula~(2) do Teor.~\ref{mf:thm:multifocal}, \textbf{nenhuma
função do campo distingue quem escreveu} --- a marca não carrega a dignidade de
quem a fez. Se se quiser uma classificação de pessoas, terá de vir de outro lado,
e o campo não a fornece. \emph{É o argumento dele, com a prova que lhe faltava.}

% ═══════════════════════════════════════════════════════════════════════════
\section{O que fica por fechar}\label{mf:aberto}

\noindent\textbf{Três coisas, e nenhuma se disfarça.}

\medskip\noindent\textbf{(1) A energia psíquica como substância.} O livro põe um
campo de energia psíquica que \emph{coexiste} com o cérebro e se transmuta nele
por \cury{janelas de transmutação}. Este documento \textbf{não a traduz e não a
usa} --- e não porque seja falsa, mas porque não tem, no livro, nenhuma
propriedade da qual se possa deduzir alguma coisa: não se diz o que se conserva,
o que a mede, nem o que a distingue de uma redescrição. \emph{Um postulado do
qual nada se segue não tem onde encaixar}, e inventar-lhe uma estrutura aqui
seria pôr na boca do autor o que ele não disse. \textbf{Não há \pt{7}.}

\medskip\noindent\textbf{(2) A natureza «virtual» do pensamento consciente.} O
livro insiste que o dialético e o antidialético são \cury{de natureza virtual},
\cury{destituídos de realidade}, e que só o essencial é real. O
Teor.~\ref{mf:thm:duas} fecha o que era \emph{estrutural} nessa distinção --- o
que mede e o que ordena --- mas \textbf{não fecha a ontologia}, e não tem como:
\emph{real} e \emph{virtual} não são predicados que a construção saiba atribuir.
Fica assinalado que o trabalho todo do livro se faz sem eles.

\medskip\noindent\textbf{(3) O salto para a clínica.} Todos os teoremas da
Parte~I são consequências dos seis postulados, e os postulados são
\textbf{modelo}, não medida. As medidas citadas são medidas de programas. O
Teor.~\ref{mf:thm:adapta} e o Cor.~\ref{mf:cor:luto} fazem uma
\textbf{previsão falsificável} --- que a resposta a uma repetição vá com $1/G$, e
que a resposta total obedeça a $H_N\le R\le N$ com repartir a maximizar --- e
essa previsão é a que teria de ser medida fora do computador para valer alguma
coisa como psicologia. \emph{Este documento não a mediu, e não deve ser lido como
se tivesse.} É, aliás, o trabalho seguinte que ele aponta.

% ═══════════════════════════════════════════════════════════════════════════
\section{O quadro, todo}\label{mf:quadro}

{\small
\begin{longtable}{@{}p{0.05\textwidth}p{0.40\textwidth}p{0.45\textwidth}@{}}
\toprule
 & a rachadura & o fecho \\
\midrule
\endfirsthead
\toprule
 & a rachadura & o fecho \\
\midrule
\endhead
\bottomrule
\multicolumn{3}{@{}r@{}}{\footnotesize\itshape continua na página seguinte}\\
\endfoot
\bottomrule
\endlastfoot
I & a objecção à matemática destrói a RPS de que o livro precisa &
\pt{1} pede só igualdade decidível; a perda é $\Phi$ --- Cor.~\ref{mf:cor:debito} \\
II & o \emph{eu} é entrada da 2.\textsuperscript{a} etapa e saída da 5.\textsuperscript{a} &
Teor.~\ref{mf:thm:zetamu}: $\mu$ e $\zeta$; Cor.~\ref{mf:cor:vitima}: um passo \\
III & o \emph{eu} distingue quem construiu, e nada regista o autor &
Teor.~\ref{mf:thm:multifocal} e Teor.~\ref{mf:thm:folha}: a ordem sim, a autoria
não \\
IV & não há lembrança, e mesmo assim há relação lógica &
Teor.~\ref{mf:thm:reconstroi} e Cor.~\ref{mf:cor:debito} \\
V & a âncora é declarada indefinível e suporta o edifício &
\pt{4} e Teor.~\ref{mf:thm:ancora}: a bola, e a partição \\
VI & o RAM reforça o que a psicoadaptação apaga &
Teor.~\ref{mf:thm:adapta} ($1/G$) e Cor.~\ref{mf:cor:luto} ($H_N\le R\le N$) \\
VII & o campo nunca repousa, e a terapia fá-lo repousar &
Teor.~\ref{mf:thm:duas}: a cisão é única, e as metades são disjuntas \\
VIII & três pensamentos, e um fenómeno tapa-buraco por passagem &
Teor.~\ref{mf:thm:duas} e Cor.~\ref{mf:cor:ponte}: a passagem não paga \\
IX & distância infinita, e mesmo assim graduável &
Teor.~\ref{mf:thm:rastro}: a distância é $\delta$, e halva \\
X & a traição é inevitável, e sem mecanismo nem custo &
Teor.~\ref{mf:thm:outro} e Cor.~\ref{mf:cor:outroeu} \\
XI & o caos é o motor, e ninguém paga a conta &
Teor.~\ref{mf:thm:custo}: $\kappa\Phi$, e o dual guardado é o piso não pago \\
XII & escravo da âncora, e capaz de a alargar &
Teor.~\ref{mf:thm:ancora}(3), e $\arg\max$ contra $\arg\min$ \\
XIII & uma só mente contra dois territórios que não se tocam &
Teor.~\ref{mf:thm:ancora}(2): uma partição não faz dois conjuntos \\
XIV & «os computadores jamais», sem critério &
memória no agente contra memória no espaço --- \code{fis:thm:mult}(3),
\code{fis:thm:periodo} \\
\end{longtable}}

\medskip\noindent\textbf{E a leitura do quadro, em uma linha.} Doze dos catorze
fechos são o \textbf{mesmo objecto} visto de ângulos diferentes: a realização
$\pi\colon I\to X$, a sua multiplicidade $G$, e o par $\zeta/\mu$ que a escreve e
a desescreve. \emph{Não é que a construção sirva para tudo}; é que o livro
descrevia, com catorze nomes, um objecto só --- e por não ter o objecto, teve de
postular um fenómeno de cada vez que precisou de outra das suas leituras.
\[
\boxed{\;\text{a disponibilidade é }\zeta;\qquad
\text{a resposta é }\mu;\qquad
\text{a adaptação é }1/G\;}
\]
\emph{Três nomes, um campo.}

\end{document}
